\documentclass{article}
\usepackage[margin=0.8in]{geometry}
\usepackage{amsmath,amssymb,amsthm,mathtools}
\usepackage{mathrsfs,bm,bbm}
\usepackage{graphicx,booktabs,array,multirow,tabularx,longtable}
\usepackage{enumitem}
\usepackage{xcolor}
\usepackage{algorithm,algpseudocode}
\usepackage{subcaption,tikz}
\usepackage{siunitx,cancel,accents}
\usepackage{microtype}
\usepackage[numbers,sort&compress]{natbib}
\usepackage[colorlinks=true,linkcolor=blue,citecolor=blue,urlcolor=blue]{hyperref}
\usepackage{cleveref}
\setlength{\emergencystretch}{2em}
\newtheorem{assumption}{Assumption}
\newtheorem{definition}{Definition}
\newtheorem{theorem}{Theorem}
\newtheorem{lemma}{Lemma}
\newtheorem{proposition}{Proposition}
\newtheorem{openendedquestion}{Open-ended Question}
\newtheorem{conjecture}{Conjecture}
\newcommand{\coreref}[1]{\ref{#1}}

\begin{document}

\sloppy

\section*{eid\_\allowbreak{}unlabeled\_\allowbreak{}perfectmix\_\allowbreak{}allcond\_\allowbreak{}certificate}

\noindent\textit{Specialization:} v1\par

\paragraph{TL;DR.} An unordered pooled mixture containing one observational centered Gaussian SEM component and one effective perfect stochastic intervention per labeled node determines its entire semantic explanation without faithfulness. Coordinate-isolated precision pairs identify baseline--root pairs, invariant root marginals orient the baseline, and acyclicity makes row-support target matching unique. These facts yield a sound-and-complete \(O(p^4)\) exact membership certificate on the observed-coordinate fully covered class, plus deterministic discrete recovery under operator error below \(\min\{\lambda/2,\alpha\lambda^2/16,g/8\}\). A verified reverse-topological relabeling of latent nodes, with \(G_{\mathrm{Sq}}=Q^{\top}\) and \(H_{\mathrm{Sq}}=Q\), shows that the legal chain, fork, and open star counterfamilies refute the ICML 2023 rank selector, and indeed every equivariant rank-only selector, on the full published class; this negative transfer does not extend the positive certificate to latent mixing. The end-to-end growing-dimensional component-estimation tail remains an explicit open question.

\section{Project justification}

\paragraph{Gap.} Appendix F.2, Proposition 10 of Squires et al. is the published unknown-baseline rank reduction, but its decisive intervention-versus-intervention rank inequality fails on legal models inside its full domain. Rank compression is more fundamentally insufficient: an embedded open star family defeats every component-permutation-equivariant rank-only selector. The accepted-bank paper \texttt{eid\_crl\_coverratio\_mmd\_genericity\_v1} instead treats separately indexed laws from faithful positive nonlinear non-Gaussian compact-support mechanisms under an arbitrary shared \(C^2\) diffeomorphic mixing map, with one unknown-target perfect intervention per latent node and ratio/MMD genericity. The present problem has one pooled unlabeled centered-Gaussian mixture, permits unfaithfulness, and directly observes the causal coordinates. Neither result subsumes the other.

\paragraph{Niche.} The observed-coordinate, centered, fully covered perfect-stochastic-intervention subclass has enough entrywise precision geometry for universal identification at every legal parameter point, including unfaithful laws and nonroot interventions that preserve marginal variance. It supports exact covariance-tuple membership and full linear-Gaussian SEM recovery. This positive niche is a strict subclass of the Squires et al. latent-mixing domain, but it is not a specialization of the nonlinear, non-Gaussian, separately indexed accepted-bank result because the observation and sampling structures differ.

\paragraph{Fill.} The coordinate certificate recovers the baseline, target bijection, weighted DAG, structural and intervention variances, and mixture weights, and it certifies nonmembership in \(O(p^4)\) exact real arithmetic. The repaired normalized-mechanism embedding uses a reverse-topological latent-node bijection \(\rho\), \(G_{\mathrm{Sq}}=Q^{\top}\), and \(H_{\mathrm{Sq}}=Q\); it preserves context labels and all context precisions. Therefore the chain, fork, and star witnesses invalidate the published rank selector on its full class, while the entrywise certificate repairs baseline selection only on the observed-coordinate fully covered subclass.

\section{Related work}

Squires et al. (2023, Theorem 1 and Proposition 1) introduce the causal-disentanglement setup and precision-difference construction. Their Appendix F.2, Proposition 10 asserts unique rank-score baseline selection on the full published class. The embedding proved here first sends the latent nodes through a reverse-topological bijection \(\rho\), then takes \(G_{\mathrm{Sq}}=Q^{\top}\), \(H_{\mathrm{Sq}}=Q\), and conjugates each normalized mechanism by \(Q\). This verifies their ordering, generic-row, and row-normalization assumptions while keeping every context label and context precision unchanged. The chain, fork, and star failures therefore refute Proposition 10 on its own domain. The coordinate certificate is a positive repair only for the embedded observed-coordinate fully covered subclass; it does not solve general latent mixing. Appendix G leaves complete membership testing open. The accepted-bank paper \texttt{eid\_crl\_coverratio\_mmd\_genericity\_v1} studies separately indexed observational and interventional laws generated by faithful positive nonlinear non-Gaussian compact-support mechanisms, allows an arbitrary shared \(C^2\) diffeomorphic mixing map, assumes one unknown-target perfect intervention per latent node, and uses ratio/MMD genericity. In contrast, this note observes one pooled unlabeled centered-Gaussian mixture, allows unfaithfulness, directly observes the causal coordinates, and gives exact covariance-tuple membership together with full linear-Gaussian SEM recovery. Neither paper subsumes the other. Kumar and Sinha (2021, Definition 3.1, Assumptions 3.1 and 3.2, and Theorem 3.1) give an exact-oracle identification result for a finite discrete causal Bayesian network when \(G\), \(P(V)\), and \(P_{\mathrm{mix}}(V)\) are supplied: under exclusion and positivity, their theorem recovers intervention tuples and weights. It neither estimates unlabeled component covariances from samples nor hides the baseline observational law. Kumar, Shiragur, and Uhler (2024, Theorem 4.1) motivate the pooled-mixture front end but fix the number of components; their Corollary 4.1.1 also uses separate observational samples and interventional faithfulness. Varici et al. (2021, Equation (6) and Theorems 1 and 3) give precision-support machinery for a designated observational/interventional pair, with faithfulness-based guarantees, but do not select a hidden baseline. Squires, Wang, and Uhler (2020, Algorithm 1 and Theorem 1) prove UT-IGSP consistency for segregated regimes under interventional faithfulness. Gamella et al. characterize \(I\)-equivalence for separately indexed linear-SCM environments with a common coefficient matrix and noise-variance interventions under faithfulness, intervention heterogeneity, and truthfulness; this is Theorem 6 in the 2022 manuscript and Theorem 1 in arXiv version 4. Their environments remain labeled, and they provide neither an unlabeled-mixture estimation front end nor perfect row-deletion surgery. Leyes Carreno, Meroni, and Seigal (2025, Definition 1.5 and Theorem 1.6) study perfect-intervention recovery with a known observational context and non-Gaussian higher-order information. Anderson et al. (2026, Theorem 1.1) and Koltchinskii and Lounici (2017, Theorems 4 and 6) supply techniques relevant to the still-open growing-component mixture-estimation tail.

\section{Setup}

Target (estimand / structural parameter): \(\theta_{\mathrm{sem}}=(c_0,\vartheta,G,B,\omega,\tau,\pi)\).
Primary functional / estimator / diagnostic: \(\mathcal R(P_{\mathrm{mix}})=\mathsf{Cert}(\{(\Sigma_c,\pi_c)\}_{c\in C})=\theta_{\mathrm{sem}}\).
\begin{itemize}
  \item \texttt{p} --- \texttt{integer}; \(\{2,3,\ldots\}\); \texttt{Number of labeled observed coordinates.}
  \item \texttt{V} --- \texttt{finite labeled vertex set}; \(V=[p]=\{1,\ldots,p\}\); \texttt{Causal vertices and intervention targets.}
  \par\smallskip\noindent\emph{Definition.} \texttt{The coordinate labels are also the causal-\allowbreak{}node labels.}
  \item \texttt{C} --- \texttt{finite component-\allowbreak{}label set}; \(C=\{0\}\cup V\); \texttt{Latent labels of the unordered Gaussian components.}
  \par\smallskip\noindent\emph{Definition.} \texttt{Truth labels for one baseline component and one component per target;\allowbreak{} the observed collection forgets this ordering.}
  \item \texttt{Ip} --- \texttt{matrix}; \(\mathbb R^{p\times p}\); \texttt{Structural-\allowbreak{}equation normalization.}
  \par\smallskip\noindent\emph{Definition.} \(I_p\) is the identity matrix.
  \item \texttt{e} --- \texttt{coordinate-\allowbreak{}basis map}; \(V\to\mathbb R^p\); \texttt{Coordinate surgery direction.}
  \par\smallskip\noindent\emph{Definition.} \(e_j\) is the vector with a one in coordinate \(j\) and zeros elsewhere.
  \item \texttt{G} --- \texttt{directed acyclic graph}; DAGs on \(V\); \texttt{Baseline causal graph.}
  \item \texttt{Pa} --- \texttt{parent-\allowbreak{}set map}; \(V\to 2^V\); \texttt{Local graph structure.}
  \par\smallskip\noindent\emph{Definition.} \(\operatorname{Pa}(j)=\{i\in V:i\to j\text{ in }G\}\).
  \item \texttt{B} --- \texttt{structural coefficient matrix}; \(\mathbb R^{p\times p}\); \texttt{Weighted adjacency matrix in equation-\allowbreak{}row convention.}
  \par\smallskip\noindent\emph{Definition.} \(B_{ji}\neq0\) if and only if \(i\in\operatorname{Pa}(j)\), and \(B_{jj}=0\).
  \item \texttt{A} --- \texttt{structural equation matrix}; \(\mathbb R^{p\times p}\); \texttt{Baseline equation operator.}
  \par\smallskip\noindent\emph{Definition.} \(A=I_p-B\).
  \item \texttt{a} --- \texttt{equation-\allowbreak{}row vector map}; \(V\to\mathbb R^p\); Transpose of baseline equation row \(j\).
  \par\smallskip\noindent\emph{Definition.} \(a_j=A^{\top}e_j\).
  \item \texttt{omega} --- \texttt{positive baseline noise-\allowbreak{}variance vector}; \((0,\infty)^p\); \texttt{Baseline structural-\allowbreak{}noise variances.}
  \item \texttt{Omega} --- \texttt{positive diagonal matrix}; \(\mathbb S_{++}^p\); \texttt{Baseline structural-\allowbreak{}noise covariance.}
  \par\smallskip\noindent\emph{Definition.} \(\Omega=\operatorname{diag}(\omega_1,\ldots,\omega_p)\).
  \item \texttt{tau} --- \texttt{positive intervention-\allowbreak{}variance vector}; \((0,\infty)^p\); \texttt{Variances of the independent stochastic intervention draws.}
  \item \texttt{Aint} --- \texttt{intervened equation-\allowbreak{}matrix map}; \(V\to\mathbb R^{p\times p}\); Perfect deletion of the incoming mechanism at target \(j\).
  \par\smallskip\noindent\emph{Definition.} \(A^{(j)}\) equals \(A\) except that row \(j\) is replaced by \(e_j^{\top}\).
  \item \texttt{Omegaint} --- \texttt{intervened noise-\allowbreak{}covariance map}; \(V\to\mathbb S_{++}^p\); Noise covariance after target \(j\) receives a perfect stochastic intervention.
  \par\smallskip\noindent\emph{Definition.} \(\Omega^{(j)}=\Omega+(\tau_j-\omega_j)e_je_j^{\top}\).
  \item \texttt{Sigma} --- \texttt{component covariance map}; \(C\to\mathbb S_{++}^p\); \texttt{Covariances of the centered Gaussian components.}
  \item \texttt{K} --- \texttt{component precision map}; \(C\to\mathbb S_{++}^p\); \texttt{Observable precision matrices.}
  \par\smallskip\noindent\emph{Definition.} \(K_c=\Sigma_c^{-1}\).
  \item \texttt{pi} --- \texttt{mixing-\allowbreak{}weight vector}; \((0,1)^{p+1}\); \texttt{Positive component weights.}
  \item \texttt{Pmix} --- \texttt{probability law}; probability measures on \(\mathbb R^p\); \texttt{Observed pooled mixture law.}
  \par\smallskip\noindent\emph{Definition.} \(P_{\mathrm{mix}}=\sum_{c\in C}\pi_c\,\mathcal N(0,\Sigma_c)\).
  \item \texttt{M} --- \texttt{model class}; collections indexed by \(p\); \texttt{Domain of the identification and certification results.}
  \par\smallskip\noindent\emph{Definition.} \(\mathcal M_p\) is the fully covered unlabeled perfect-surgery class defined below.
  \item \texttt{c0} --- \texttt{component label}; \(C\); \texttt{Recovered observational baseline position in an arbitrary input ordering.}
  \item \texttt{vartheta} --- \texttt{bijection}; \(C\setminus\{c_0\}\to V\); \texttt{Recovered assignment of nonbaseline components to intervention targets.}
  \item \texttt{thetaSem} --- \texttt{semantic parameter}; \texttt{baseline label,\allowbreak{} target bijection,\allowbreak{} DAG,\allowbreak{} coefficients,\allowbreak{} variances,\allowbreak{} and weights}; \texttt{Causal structure and parameter estimand.}
  \par\smallskip\noindent\emph{Definition.} \(\theta_{\mathrm{sem}}=(c_0,\vartheta,G,B,\omega,\tau,\pi)\).
  \item \texttt{Cert} --- \texttt{partial deterministic map}; unordered covariance-weight tuples \(\to\) semantic explanations or rejection traces; \texttt{Sound-\allowbreak{}and-\allowbreak{}complete membership and recovery algorithm.}
  \par\smallskip\noindent\emph{Definition.} \(\mathsf{Cert}\) is the coordinate certificate defined below.
  \item \texttt{Recover} --- \texttt{identifying functional}; \(P_{\mathrm{mix}}\to\theta_{\mathrm{sem}}\); \texttt{Population recovery map.}
  \par\smallskip\noindent\emph{Definition.} \(\mathcal R\) first identifies the unordered finite Gaussian mixture and then applies \(\mathsf{Cert}\).
  \item \texttt{Delta} --- \texttt{precision-\allowbreak{}difference map}; \(V\to\mathbb S^p\); \texttt{Local intervention signature.}
  \par\smallskip\noindent\emph{Definition.} \(\Delta_j=K_j-K_0\) in a truth-labeled representation.
  \item \texttt{rowsupp} --- \texttt{row-\allowbreak{}support functional}; \(\mathbb R^{p\times p}\to2^V\); \texttt{Support statistic used for target matching.}
  \par\smallskip\noindent\emph{Definition.} \(\operatorname{rowsupp}(M)=\{i\in V: M_{ik}\neq0\text{ for some }k\in V\}\).
  \item \texttt{S} --- \texttt{component support map}; \(C\setminus\{c_0\}\to2^V\); \texttt{Supports supplied to singleton target matching.}
  \par\smallskip\noindent\emph{Definition.} \(S_c=\operatorname{rowsupp}(K_c-K_{c_0})\).
  \item \texttt{Roots} --- \texttt{vertex subset}; \(2^V\); \texttt{Root coordinates used to orient the baseline.}
  \par\smallskip\noindent\emph{Definition.} \(\operatorname{Roots}(G)=\{r\in V:\operatorname{Pa}(r)=\varnothing\}\).
  \item \texttt{lambda} --- \texttt{positive scalar}; \texttt{All-\allowbreak{}component conditioning margin.}
  \par\smallskip\noindent\emph{Definition.} \(\lambda=\min_{c\in C}\lambda_{\min}(\Sigma_c)\).
  \item \texttt{alpha} --- \texttt{positive scalar}; \texttt{Full-\allowbreak{}entry precision-\allowbreak{}support margin.}
  \par\smallskip\noindent\emph{Definition.} \(\alpha=\min\{|(\Delta_j)_{ik}|:j,i,k\in V,\ (\Delta_j)_{ik}\neq0\}\).
  \item \texttt{g} --- \texttt{positive scalar}; \texttt{Root marginal-\allowbreak{}variance orientation gap.}
  \par\smallskip\noindent\emph{Definition.} \(g=\min_{r\in\operatorname{Roots}(G)}|\tau_r-\omega_r|\).
  \item \texttt{epsilonStar} --- \texttt{positive scalar}; \texttt{Certified covariance-\allowbreak{}operator perturbation radius.}
  \par\smallskip\noindent\emph{Definition.} \(\varepsilon_{\star}=\min\{\lambda/2,\alpha\lambda^2/16,g/8\}\).
  \item \texttt{SigmaHat} --- \texttt{estimated covariance map}; \(C\to\mathbb S^p\); \texttt{Unordered estimated component covariances.}
  \item \texttt{KHat} --- \texttt{estimated precision map}; \(C\to\mathbb S_{++}^p\); \texttt{Precision matrices used by the stable decoder.}
  \par\smallskip\noindent\emph{Definition.} \(\widehat K_c=\widehat\Sigma_c^{-1}\) whenever the inverse exists.
  \item \texttt{Eop} --- \texttt{permutation-\allowbreak{}matched operator error}; \texttt{Observable-\allowbreak{}estimator error interface to deterministic recovery.}
  \par\smallskip\noindent\emph{Definition.} \(E_{\mathrm{op}}=\min_{\rho\in\mathfrak S_{p+1}}\max_{c\in C}\|\widehat\Sigma_{\rho(c)}-\Sigma_c\|_{\mathrm{op}}\).
  \item \texttt{Decoder} --- \texttt{deterministic decoder}; estimated covariance tuples \(\to(c_0,\vartheta,G)\); \texttt{Stable discrete recovery procedure.}
  \par\smallskip\noindent\emph{Definition.} \(\widehat{\mathsf R}_{\alpha,g}\) is the thresholded coordinate decoder defined below.
  \item \texttt{R} --- \texttt{symmetric integer matrix}; \(\{0,1,\ldots,p\}^{(p+1)\times(p+1)}\); \texttt{Pairwise rank compression of the covariance collection.}
  \par\smallskip\noindent\emph{Definition.} \(R_{cd}=\operatorname{rank}(K_c-K_d)\).
  \item \texttt{rscore} --- \texttt{rank-\allowbreak{}score vector}; \texttt{Published baseline-\allowbreak{}selection statistic and its symmetric repair.}
  \par\smallskip\noindent\emph{Definition.} \(r_c=\sum_{d\in C\setminus\{c\}}R_{cd}\), with the literal ICML score obtained by its published restricted summation convention.
  \item \texttt{n} --- \texttt{positive integer}; \texttt{Pooled sample size for the open statistical extension.}
  \item \texttt{Zsample} --- \texttt{random sample}; \((\mathbb R^p)^n\); \texttt{Observed pooled data without regime labels.}
  \par\smallskip\noindent\emph{Definition.} \(Z_1,\ldots,Z_n\stackrel{\mathrm{iid}}{\sim}P_{\mathrm{mix}}\).
  \item \texttt{delta} --- \texttt{confidence level}; \((0,1)\); \texttt{Tail-\allowbreak{}probability index for the open statistical extension.}
  \item \texttt{Tn} --- \texttt{nonnegative tail-\allowbreak{}radius functional}; \texttt{Target permutation-\allowbreak{}matched covariance operator-\allowbreak{}error radius.}
  \par\smallskip\noindent\emph{Definition.} \(T_n(\delta)\) is derived by Open-ended Question 1.
  \item \texttt{Hmix} --- \texttt{estimation handle}; pooled samples \(\to\) unordered covariance estimates; \texttt{Concrete starting move for Open-\allowbreak{}ended Question 1.}
  \par\smallskip\noindent\emph{Definition.} \(\mathsf H_{\mathrm{mix}}\) is the sum-of-squares clustering and effective-rank covariance handle defined below.
\end{itemize}

\section{Assumptions}

\begin{assumption}[ass:baseline-gaussian-sem]\label{ass:baseline-gaussian-sem}
The baseline covariance satisfies \(\Sigma_0=A^{-1}\Omega A^{-\top}\).
\par\smallskip\noindent\emph{Standard} (Acyclic linear Gaussian structural equation model, \cite{pmlr-v202-squires23a}).
\end{assumption}

\begin{assumption}[ass:perfect-single-node-surgeries]\label{ass:perfect-single-node-surgeries}
For every \(j\in V\), \(\Sigma_j=(A^{(j)})^{-1}\Omega^{(j)}(A^{(j)})^{-\top}\).
\par\smallskip\noindent\emph{Standard} (One stochastic perfect single-node intervention per node, \cite{pmlr-v202-squires23a}).
\end{assumption}

\begin{assumption}[ass:effective-interventions]\label{ass:effective-interventions}
For every \(j\in V\), \(\operatorname{Pa}(j)\neq\varnothing\) or \(\tau_j\neq\omega_j\).
\par\smallskip\noindent\emph{Standard} (Effective-intervention condition specialized to centered perfect interventions, \cite{kumarShiragurUhler2024mixtures}).
\end{assumption}

\begin{assumption}[ass:positive-mixture-weights]\label{ass:positive-mixture-weights}
The weights satisfy \(\pi_c>0\) for every \(c\in C\) and \(\sum_{c\in C}\pi_c=1\).
\par\smallskip\noindent\emph{Standard} (Positive finite-mixture weights, \cite{kumarShiragurUhler2024mixtures}).
\end{assumption}

\begin{assumption}[ass:iid-pooled-sample]\label{ass:iid-pooled-sample}
The observations \(Z_1,\ldots,Z_n\) are independent draws from \(P_{\mathrm{mix}}\).
\par\smallskip\noindent\emph{Standard} (Independent pooled-mixture sampling, \cite{kumarShiragurUhler2024mixtures}).
\end{assumption}

\section{Definitions}

\begin{definition}[def:model-class]\label{def:model-class}
\par\smallskip\noindent\emph{Defined object.} Fully covered unlabeled perfect-surgery class \(\mathcal M_p\)
\par\smallskip\noindent\emph{Construction.} \(\mathcal M_p=\{(G,B,\omega,\tau,\{\Sigma_c,\pi_c\}_{c\in C}):\Sigma_0=A^{-1}\Omega A^{-\top},\ \Sigma_j=(A^{(j)})^{-1}\Omega^{(j)}(A^{(j)})^{-\top}\ \forall j\in V,\ [\operatorname{Pa}(j)\neq\varnothing\ \text{or}\ \tau_j\neq\omega_j]\ \forall j\in V,\ \pi_c>0\ \forall c\in C,\ \sum_{c\in C}\pi_c=1\}\).
\par\smallskip\noindent\emph{Member properties.} \texttt{ass:\allowbreak{}baseline-\allowbreak{}gaussian-\allowbreak{}sem}; \texttt{ass:\allowbreak{}perfect-\allowbreak{}single-\allowbreak{}node-\allowbreak{}surgeries}; \texttt{ass:\allowbreak{}effective-\allowbreak{}interventions}; \texttt{ass:\allowbreak{}positive-\allowbreak{}mixture-\allowbreak{}weights}.
\end{definition}

\begin{definition}[def:mixture-law]\label{def:mixture-law}
\par\smallskip\noindent\emph{Defined object.} \texttt{Pooled law}
\par\smallskip\noindent\emph{Construction.} \(P_{\mathrm{mix}}=\sum_{c\in C}\pi_c\mathcal N(0,\Sigma_c)\).
\par\smallskip\noindent\emph{Inputs.} \texttt{Sigma}; \texttt{pi}.
\end{definition}

\begin{definition}[def:precision-differences]\label{def:precision-differences}
\par\smallskip\noindent\emph{Defined object.} \texttt{Baseline precision differences}
\par\smallskip\noindent\emph{Construction.} \(\Delta_j=K_j-K_0\) for every \(j\in V\).
\par\smallskip\noindent\emph{Inputs.} \texttt{K}.
\end{definition}

\begin{definition}[def:row-support]\label{def:row-support}
\par\smallskip\noindent\emph{Defined object.} \texttt{Row support}
\par\smallskip\noindent\emph{Construction.} \(\operatorname{rowsupp}(M)=\{i\in V: M_{ik}\neq0\text{ for some }k\in V\}\).
\par\smallskip\noindent\emph{Inputs.} \texttt{V}.
\end{definition}

\begin{definition}[def:coordinate-isolated-pair]\label{def:coordinate-isolated-pair}
\par\smallskip\noindent\emph{Defined object.} \texttt{Coordinate-\allowbreak{}isolated precision pair}
\par\smallskip\noindent\emph{Construction.} A pair \(\{c,d\}\subset C\) is coordinate-isolated at \(r\in V\) exactly when \(K_c-K_d=s e_re_r^{\top}\) for some \(s\neq0\).
\par\smallskip\noindent\emph{Inputs.} \texttt{K}.
\end{definition}

\begin{definition}[def:semantic-estimand]\label{def:semantic-estimand}
\par\smallskip\noindent\emph{Defined object.} \texttt{Semantic explanation}
\par\smallskip\noindent\emph{Construction.} \(\theta_{\mathrm{sem}}=(c_0,\vartheta,G,B,\omega,\tau,\pi)\), where \(c_0\) is an input position and \(\vartheta:C\setminus\{c_0\}\to V\) is a bijection.
\par\smallskip\noindent\emph{Inputs.} \texttt{Sigma}; \texttt{pi}.
\end{definition}

\begin{definition}[def:exact-certificate]\label{def:exact-certificate}
\par\smallskip\noindent\emph{Defined object.} Coordinate certificate \(\mathsf{Cert}\)
\par\smallskip\noindent\emph{Construction.} Given an unordered \(p+1\)-tuple \(\{(\Sigma_c,\pi_c)\}\), check positive definiteness, distinctness, and weights; invert every covariance; find and orient a coordinate-isolated pair using a third component's marginal; form every support \(\operatorname{rowsupp}(K_c-K_{c_0})\); repeatedly match a component whose residual support contains one unmatched coordinate; recover each row by baseline regression; then verify \(K_{c_0}=A^{\top}\Omega^{-1}A\) and, for every matched pair \(\vartheta(c)=j\), verify \(K_c-K_{c_0}=\tau_j^{-1}e_je_j^{\top}-\omega_j^{-1}a_ja_j^{\top}\).
\par\smallskip\noindent\emph{Inputs.} \texttt{Sigma}; \texttt{pi}.
\end{definition}

\begin{definition}[def:stability-margins]\label{def:stability-margins}
\par\smallskip\noindent\emph{Defined object.} \texttt{All-\allowbreak{}component recovery margins}
\par\smallskip\noindent\emph{Construction.} \(\lambda=\min_{c\in C}\lambda_{\min}(\Sigma_c)\), \(\alpha=\min\{|(\Delta_j)_{ik}|:(\Delta_j)_{ik}\neq0\}\), \(g=\min_{r\in\operatorname{Roots}(G)}|\tau_r-\omega_r|\), and \(\varepsilon_{\star}=\min\{\lambda/2,\alpha\lambda^2/16,g/8\}\).
\par\smallskip\noindent\emph{Inputs.} \texttt{Sigma}; \texttt{Delta}; \texttt{G}; \texttt{tau}; \texttt{omega}.
\end{definition}

\begin{definition}[def:threshold-decoder]\label{def:threshold-decoder}
\par\smallskip\noindent\emph{Defined object.} Thresholded coordinate decoder \(\widehat{\mathsf R}_{\alpha,g}\)
\par\smallskip\noindent\emph{Construction.} Invert \(\widehat\Sigma_c\); detect an entry of \(\widehat K_c-\widehat K_d\) when its magnitude exceeds \(\alpha/2\); recognize a coordinate pair by one detected diagonal and no other detected entry; orient it toward the endpoint whose corresponding marginal is closer to any third component; then apply singleton row-support matching.
\par\smallskip\noindent\emph{Inputs.} \texttt{SigmaHat}; \texttt{alpha}.
\end{definition}

\begin{definition}[def:rank-compression]\label{def:rank-compression}
\par\smallskip\noindent\emph{Defined object.} \texttt{Rank-\allowbreak{}only baseline information}
\par\smallskip\noindent\emph{Construction.} \(R_{cd}=\operatorname{rank}(K_c-K_d)\) and \(r_c=\sum_{d\neq c}R_{cd}\); a rank-only selector is a component-permutation-equivariant map from \(R\) to one vertex of \(C\).
\par\smallskip\noindent\emph{Inputs.} \texttt{K}.
\end{definition}

\begin{definition}[def:two-node-counterfamilies]\label{def:two-node-counterfamilies}
\par\smallskip\noindent\emph{Defined object.} \texttt{Two-\allowbreak{}node chain counterfamilies}
\par\smallskip\noindent\emph{Construction.} For \(1\to2\), set \(B_{21}=b\neq0\), \(\omega=(u,v)\), and \(\tau=(t,s)\), so \(\Sigma_0=\begin{pmatrix}u&bu\\bu&b^2u+v\end{pmatrix}\), \(\Sigma_1=\begin{pmatrix}t&bt\\bt&b^2t+v\end{pmatrix}\), and \(\Sigma_2=\operatorname{diag}(u,s)\); use \((b,u,v,t,s)=(2,1,1,3,2)\) and \((1,1,1,1/2,2)\) as legal witnesses.
\par\smallskip\noindent\emph{Inputs.} \texttt{B}; \texttt{omega}; \texttt{tau}.
\end{definition}

\begin{definition}[def:star-counterfamily]\label{def:star-counterfamily}
\par\smallskip\noindent\emph{Defined object.} \texttt{One-\allowbreak{}root star rank-\allowbreak{}obstruction family}
\par\smallskip\noindent\emph{Construction.} For every \(p\ge2\), take the one-root star with root \(1\), edges \(1\to j\) of coefficients \(b_j\neq0\) for \(j=2,\ldots,p\), positive baseline variances \(\omega_j\), and positive intervention variances \(\tau_j\). Put \(d=\tau_1^{-1}-\omega_1^{-1}\neq0\) and restrict to the nonempty strict region \(d(\tau_j-\omega_j)-b_j^2<0\) for every leaf \(j\).
\par\smallskip\noindent\emph{Inputs.} \texttt{p}.
\end{definition}

\begin{definition}[def:three-node-counterfamily]\label{def:three-node-counterfamily}
\par\smallskip\noindent\emph{Defined object.} \texttt{Three-\allowbreak{}node fork counterfamily}
\par\smallskip\noindent\emph{Construction.} Take the fork \(1\to2\), \(1\to3\) with coefficients \(2\) and \(3\), \(\omega=(1,1,1)\), and \(\tau=(3,4,5)\); its pairwise rank matrix is \(\begin{pmatrix}0&1&2&2\\1&0&2&2\\2&2&0&3\\2&2&3&0\end{pmatrix}\). Retain the open parameter neighborhood on which these ranks remain unchanged.
\par\smallskip\noindent\emph{Inputs.} \texttt{G}; \texttt{B}; \texttt{omega}; \texttt{tau}.
\end{definition}

\begin{definition}[def:false-positive-tuple]\label{def:false-positive-tuple}
\par\smallskip\noindent\emph{Defined object.} \texttt{Local-\allowbreak{}surgery false-\allowbreak{}positive tuple}
\par\smallskip\noindent\emph{Construction.} In dimension three, let \(K_0=\begin{pmatrix}1&0&0\\0&2&-1\\0&-1&2\end{pmatrix}\) and \(K_j=K_0+e_je_j^{\top}\) for \(j=1,2,3\).
\end{definition}

\begin{definition}[def:near-singular-family]\label{def:near-singular-family}
\par\smallskip\noindent\emph{Defined object.} \texttt{Near-\allowbreak{}singular all-\allowbreak{}component conditioning family}
\par\smallskip\noindent\emph{Construction.} For \(0<t<1/3\), let \(\Sigma_0=\begin{pmatrix}1&1\\1&2\end{pmatrix}\), \(\Sigma_1(t)=\begin{pmatrix}t&t\\t&1+t\end{pmatrix}\), \(\Sigma_2=\operatorname{diag}(1,2)\), and perturb only component one to \(\widehat\Sigma_1(t)=\begin{pmatrix}t&2t\\2t&1+t\end{pmatrix}\).
\end{definition}

\begin{definition}[def:mixture-estimation-handle]\label{def:mixture-estimation-handle}
\par\smallskip\noindent\emph{Defined object.} \texttt{Growing-\allowbreak{}component estimation handle}
\par\smallskip\noindent\emph{Construction.} Apply a separation-preserving sum-of-squares projection and recursive partial clustering to the centered pooled sample, estimate each recovered cluster covariance with effective-rank control, and permutation-match the resulting \(p+1\) covariance estimates before passing them to the coordinate certificate.
\par\smallskip\noindent\emph{Inputs.} \texttt{Zsample}.
\end{definition}

\section{Main results}

\begin{lemma}[lem:component-distinctness]\label{lem:component-distinctness}
Every member of \(\mathcal M_p\) has \(p+1\) distinct covariance matrices, and the pooled law \(P_{\mathrm{mix}}\) identifies the unordered covariance-weight pairs \(\{(\Sigma_c,\pi_c):c\in C\}\) up to permutation.
\end{lemma}
\begin{proof}
Fix a member of \(\mathcal M_p\).  Acyclicity supplies a topological rank
\(q:V\to\{1,\ldots,p\}\), strictly increasing along every directed edge.
The unit-diagonal matrix \(A=I_p-B\) is triangular after the corresponding
simultaneous permutation of rows and columns, and hence it and every
\(A^{(j)}\) are invertible.  Inverting the covariance factorizations in
\(\texttt{def:model-class}\) gives
\[
 K_0=A^{\top}\Omega^{-1}A
     =\sum_{\ell\in V}\omega_\ell^{-1}a_\ell a_\ell^{\top}
\]
and, because surgery changes only equation row \(j\) and its noise variance,
\[
 K_j-K_0
 =\tau_j^{-1}e_je_j^{\top}
  -\omega_j^{-1}a_ja_j^{\top}.                                      \tag{1}
\]
Here
\[
 a_j=e_j-\sum_{i\in\operatorname{Pa}(j)}B_{ji}e_i.                   \tag{2}
\]
If \(j\) has a parent \(i\), the \((j,i)\) entry of (1) is
\(B_{ji}/\omega_j\neq0\).  If \(j\) is a root, (1) is
\[
 (\tau_j^{-1}-\omega_j^{-1})e_je_j^{\top}\neq0
\]
by \(\texttt{ass:effective-interventions}\).  Thus \(K_j\neq K_0\) for
every \(j\).

For distinct \(j,k\), relabel them so that \(q(j)<q(k)\).  Equation (2)
shows that row \(k\) of \(K_j-K_0\) is zero: \(k\) is neither \(j\) nor a
parent of \(j\).  Row \(k\) of \(K_k-K_0\) is nonzero, by the preceding
parent-entry argument when \(k\) is nonroot and by the effective-root
diagonal argument when \(k\) is a root.  Hence \(K_j\neq K_k\).  Matrix
inversion is injective on \(\mathbb S_{++}^p\), so the \(p+1\) covariance
matrices are pairwise distinct.

It remains to pass from the pooled law to its unordered component tuple.
We prove the required finite-mixture fact directly.  Suppose two finite
representations by centered nonsingular Gaussian laws have the same
probability law.  Subtract their characteristic functions and collect
identical covariance matrices.  This yields distinct matrices
\(C_1,\ldots,C_m\in\mathbb S_{++}^p\) and real coefficients
\(h_1,\ldots,h_m\) such that
\[
 \sum_{\ell=1}^m h_\ell
 \exp\!\left(-\frac12 t^{\top}C_\ell t\right)=0
 \qquad\text{for every }t\in\mathbb R^p.                            \tag{3}
\]
For each \(\ell\neq r\), the polynomial
\(v\mapsto v^{\top}(C_\ell-C_r)v\) is nonzero.  Their finite product is a
nonzero real polynomial, so there is a \(v\in\mathbb R^p\) outside its
zero set.  Consequently the positive numbers
\(s_\ell=v^{\top}C_\ell v\) are pairwise distinct.  Substituting
\(t=xv\) into (3), ordering \(s_1<\cdots<s_m\), and multiplying by
\(\exp(s_1x^2/2)\) gives
\[
 h_1+\sum_{\ell=2}^m h_\ell
       \exp\!\left(-\frac12(s_\ell-s_1)x^2\right)=0 .
\]
Letting \(x\to\infty\) proves \(h_1=0\).  Iteration proves every
\(h_\ell=0\).  Thus the discrete mixing measure on covariance matrices is
unique.  In particular, the positive weights required by
\(\texttt{ass:positive-mixture-weights}\) and their associated distinct
covariances in \(\texttt{def:mixture-law}\) are determined up to one common
permutation.
\end{proof}
\par\smallskip\noindent \textit{Justification.} This lemma closes the population passage from a pooled distribution to the unordered component tuple. \textit{Gap.} Kumar, Shiragur, and Uhler identify general mixture parameters for fixed component count; this lemma isolates the exact population fact required when the causal labels are still hidden. \textit{Consumer.} Theorem 1 can start from the pooled law rather than assuming that component covariances arrive with semantic labels.

\begin{lemma}[lem:surgery-identity]\label{lem:surgery-identity}
For every \(j\in V\), writing \(a_j=A^{\top}e_j\), the perfect surgery obeys \(\Delta_j=\tau_j^{-1}e_je_j^{\top}-\omega_j^{-1}a_ja_j^{\top}\) and \(\operatorname{rowsupp}(\Delta_j)=\{j\}\cup\operatorname{Pa}(j)\), including when a nonroot has \(\tau_j=\omega_j\).
\end{lemma}
\begin{proof}
By \(\texttt{ass:baseline-gaussian-sem}\), acyclicity makes \(A\)
invertible and
\[
 K_0=\Sigma_0^{-1}=A^{\top}\Omega^{-1}A
     =\sum_{\ell\in V}\omega_\ell^{-1}a_\ell a_\ell^{\top}.          \tag{4}
\]
By \(\texttt{ass:perfect-single-node-surgeries}\), \(A^{(j)}\) has the
same rows as \(A\) except that row \(j\) is \(e_j^{\top}\), while the
corresponding diagonal noise variance is \(\tau_j\).  Therefore
\[
 K_j=(A^{(j)})^{\top}(\Omega^{(j)})^{-1}A^{(j)}
 =\sum_{\ell\neq j}\omega_\ell^{-1}a_\ell a_\ell^{\top}
  +\tau_j^{-1}e_je_j^{\top}.
\]
Subtracting (4) proves
\[
 \Delta_j
 =\tau_j^{-1}e_je_j^{\top}
  -\omega_j^{-1}a_ja_j^{\top}.                                    \tag{5}
\]

The equation-row convention gives
\[
 a_j=e_j-\sum_{i\in\operatorname{Pa}(j)}B_{ji}e_i.
\]
Thus no row outside \(\{j\}\cup\operatorname{Pa}(j)\) can be nonzero in
(5).  For every \(i\in\operatorname{Pa}(j)\),
\[
 (\Delta_j)_{ji}=(\Delta_j)_{ij}=\frac{B_{ji}}{\omega_j}\neq0,       \tag{6}
\]
so every parent row and the target row are nonzero.  If the parent set is
empty, (5) reduces to
\[
 \Delta_j=(\tau_j^{-1}-\omega_j^{-1})e_je_j^{\top},
\]
whose target row is nonzero by
\(\texttt{ass:effective-interventions}\).  These cases prove
\[
 \operatorname{rowsupp}(\Delta_j)
 =\{j\}\cup\operatorname{Pa}(j).
\]
The argument for a nonroot used (6), not
\(\tau_j\neq\omega_j\), so it also covers \(\tau_j=\omega_j\).
\end{proof}
\par\smallskip\noindent \textit{Justification.} The exact row support is the information discarded by the published rank compression. \textit{Gap.} Squires et al. give a rank-two difference identity, while Varici et al. exploit sparse precision differences for labeled pairs; neither states this universal unlabeled perfect-surgery row-support equality with unchanged nonroot variance allowed. \textit{Consumer.} The row supports drive baseline recognition, unique target matching, and the stability threshold.

\begin{lemma}[lem:coordinate-pair-characterization]\label{lem:coordinate-pair-characterization}
A pair \(\{c,d\}\subset C\) is coordinate-isolated at \(r\in V\) if and only if its endpoints are the baseline and the perfect intervention on a root \(r\); for any third component, exactly the baseline endpoint has the same \(r\)-th marginal variance as that third component.
\end{lemma}
\begin{proof}
Suppose first that one endpoint is the baseline and the other is target
\(j\).  By \(\texttt{lem:surgery-identity}\), their precision difference
has row support \(\{j\}\cup\operatorname{Pa}(j)\).  It equals a nonzero
multiple of \(e_re_r^{\top}\) exactly when this support is the singleton
\(\{r\}\), which is equivalent to \(r=j\) and
\(\operatorname{Pa}(j)=\varnothing\).  Effectiveness guarantees that the
multiple is nonzero.  Hence every baseline--root pair, and no other pair
having a baseline endpoint, is coordinate-isolated.

It remains to exclude a pair of intervention components.  Let their
distinct targets be \(u,v\), ordered so that \(u\) precedes \(v\) in a
topological order.  Write \(\Delta_\ell=K_\ell-K_0\).  Row \(v\) of
\(\Delta_u\) is zero because \(v\notin\{u\}\cup\operatorname{Pa}(u)\),
whereas row \(v\) of \(\Delta_v\) is nonzero by
\(\texttt{lem:surgery-identity}\).  If
\[
 K_v-K_u=\Delta_v-\Delta_u=s e_re_r^{\top},\qquad s\neq0,            \tag{7}
\]
then (7) forces \(r=v\).  All off-diagonal entries in row \(v\) of
\(\Delta_v\) must then vanish.  Formula (6) in the proof of
\(\texttt{lem:surgery-identity}\) forces
\(\operatorname{Pa}(v)=\varnothing\), so \(\Delta_v\) is supported only
at \((v,v)\).  But row \(u\) of \(\Delta_u\) is nonzero and row \(u\) of
\(\Delta_v\) is zero; hence row \(u\) of \(\Delta_v-\Delta_u\) is
nonzero, contradicting (7).  This proves the pair characterization.

For its orientation, let \(r\) be the characterized root.  In the
baseline SEM its equation is \(X_r=\varepsilon_r\).  Every intervention
other than the one at \(r\) retains this equation and its noise variance,
so
\[
 (\Sigma_0)_{rr}=(\Sigma_z)_{rr}=\omega_r
 \quad\text{for every }z\in C\setminus\{0,r\},
 \qquad
 (\Sigma_r)_{rr}=\tau_r.
\]
The effective-root condition gives \(\tau_r\neq\omega_r\).  Since
\(p\ge2\), the set \(C\setminus\{0,r\}\) is nonempty.  Therefore, for
every third component, exactly the baseline endpoint has the same
\(r\)-th marginal variance.
\end{proof}
\par\smallskip\noindent \textit{Justification.} This is the permutation-equivariant replacement for minimizing pairwise rank sums. \textit{Gap.} Proposition 10 of Squires et al. attempts baseline recovery from ranks; no cited result characterizes baseline-root pairs by exact coordinate support and orients them by an invariant marginal. \textit{Consumer.} Theorem 1 obtains the baseline without faithfulness, a known root, or an externally supplied observational sample.

\begin{theorem}[thm:universal-recovery]\label{thm:universal-recovery}
For every member of \(\mathcal M_p\), the pooled law \(P_{\mathrm{mix}}\) uniquely determines \(\theta_{\mathrm{sem}}\). After identifying \(c_0\), each nonbaseline component \(c\) has support \(S_c=\operatorname{rowsupp}(K_c-K_{c_0})\); the unique bijection satisfies \(\vartheta(c)\in S_c\), and then \(\operatorname{Pa}(j)=S_{\vartheta^{-1}(j)}\setminus\{j\}\), \(B_{j,\operatorname{Pa}(j)}^{\top}=(\Sigma_{c_0})_{\operatorname{Pa}(j),\operatorname{Pa}(j)}^{-1}(\Sigma_{c_0})_{\operatorname{Pa}(j),j}\), \(\omega_j=(\Sigma_{c_0})_{jj}-(\Sigma_{c_0})_{j,\operatorname{Pa}(j)}(\Sigma_{c_0})_{\operatorname{Pa}(j),\operatorname{Pa}(j)}^{-1}(\Sigma_{c_0})_{\operatorname{Pa}(j),j}\), and \(\tau_j=(\Sigma_{\vartheta^{-1}(j)})_{jj}\).
\end{theorem}
\begin{proof}
Fix an arbitrary member of \(\mathcal M_p\).  Thus \(p\geq 2\), every
\(\Sigma_c\) is positive definite, the weights are positive and sum to one,
the baseline factorization in \(\texttt{ass:baseline-gaussian-sem}\) holds,
and every intervention satisfies
\(\texttt{ass:perfect-single-node-surgeries}\) and
\(\texttt{ass:effective-interventions}\).  By
\(\texttt{lem:component-distinctness}\), \(P_{\mathrm{mix}}\) determines the
unordered collection of distinct covariance--weight pairs
\[
 \bigl\{(\Sigma_c,\pi_c):c\in C\bigr\}
\]
up to a simultaneous permutation of their component positions.  In
particular, inversion is legal and determines every
\(K_c=\Sigma_c^{-1}\).  We now construct the semantic explanation from this
observable collection and prove that the construction returns the generating
one.

First we recover the baseline position.  Every finite DAG has a root: if each
vertex had a parent, repeatedly choosing a parent would revisit a vertex and
produce a directed cycle.  Let \(r\) be a root.  Then
\(a_r=e_r\), and effectiveness gives \(\tau_r\neq\omega_r\).  Hence the
surgery identity gives
\[
 K_r-K_0
 =\left(\tau_r^{-1}-\omega_r^{-1}\right)e_re_r^{\top},
 \qquad \tau_r^{-1}-\omega_r^{-1}\neq0.                 \tag{1}
\]
Thus at least one coordinate-isolated pair is present in the observed
collection.  Conversely, by
\(\texttt{lem:coordinate-pair-characterization}\), every pair isolated at a
coordinate \(r\) consists exactly of the baseline and the intervention on the
root \(r\).  Since \(p\geq2\), the collection has \(p+1\geq3\) components, so
there is a third component.  The same lemma states that, of the two endpoints
of the isolated pair, exactly the baseline endpoint has its \(r\)-th marginal
variance equal to that of any third component.  All quantities in this test
are entries of observed covariance matrices.  It therefore orients every
coordinate-isolated pair toward one and the same position, namely the true
baseline position \(c_0\).  This proves existence and uniqueness of the
recovered \(c_0\).

For each remaining component form the observable set
\[
 S_c=\operatorname{rowsupp}(K_c-K_{c_0}),
 \qquad c\in C\setminus\{c_0\}.                         \tag{2}
\]
Temporarily use the true target label \(j\) for the component generated by
the intervention on \(j\).  The support conclusion of
\(\texttt{lem:surgery-identity}\), which also covers the allowed nonroot case
\(\tau_j=\omega_j\), gives
\[
 S_j=\{j\}\cup\operatorname{Pa}(j).                    \tag{3}
\]
Consequently the true component-to-target map is a bijection whose assigned
coordinate lies in the corresponding set in (2).

That support-compatible bijection is unique.  Indeed, let another such
bijection be given and, after identifying components by their true targets,
write it as a permutation \(f:V\to V\).  Equation (3) implies
\[
 f(j)\in\{j\}\cup\operatorname{Pa}(j)
 \quad\text{for every }j\in V.                          \tag{4}
\]
A finite acyclic graph admits a rank \(q:V\to\{1,\ldots,p\}\) that is
strictly increasing along each edge: recursively remove a root and assign
successive ranks.  Therefore (4) yields
\[
 q(f(j))\leq q(j),
\]
with strict inequality whenever \(f(j)\neq j\).  But bijectivity of \(f\)
also yields
\[
 \sum_{j\in V}q(f(j))=\sum_{j\in V}q(j).                \tag{5}
\]
Equations (4)--(5) rule out every strict inequality, and hence
\(f(j)=j\) for all \(j\).  Thus the unique observable support-compatible
matching is the true bijection \(\vartheta\).  It is a finite recovery map,
for example by enumerating the finitely many perfect matchings and retaining
the unique one satisfying \(\vartheta(c)\in S_c\).  Equation (3) now gives
the explicit graph recovery formula
\[
 \operatorname{Pa}(j)
 =S_{\vartheta^{-1}(j)}\setminus\{j\}.                  \tag{6}
\]

It remains to recover the numerical SEM parameters.  Put
\(P=\operatorname{Pa}(j)\), and let \(X\) be centered Gaussian with covariance
\(\Sigma_{c_0}\).  Define the proof intermediate \(\varepsilon=AX\).  The
baseline factorization gives, equation by equation,
\[
 \operatorname{Cov}(\varepsilon)
 =A\Sigma_{c_0}A^{\top}
 =A(A^{-1}\Omega A^{-\top})A^{\top}
 =\Omega.                                               \tag{7}
\]
Acyclicity makes \(B\) nilpotent and therefore
\[
 A^{-1}=(I_p-B)^{-1}=\sum_{m=0}^{p-1}B^m.               \tag{8}
\]
If \(i\in P\), then \(i\to j\).  For \(m\geq1\), a nonzero summand
contributing to \((B^m)_{ij}\) represents a directed path from \(j\) to
\(i\), which together with \(i\to j\) would be a directed cycle.  Also
\((I_p)_{ij}=0\).  Hence (8) implies
\((A^{-1})_{ij}=0\), and diagonality of \(\Omega\) gives
\[
 \operatorname{Cov}(X_i,\varepsilon_j)
 =e_i^{\top}A^{-1}\Omega e_j
 =\omega_j(A^{-1})_{ij}=0,
 \qquad i\in P.                                        \tag{9}
\]
The \(j\)-th row of \(\varepsilon=AX\) is
\[
 X_j=B_{jP}X_P+\varepsilon_j.                           \tag{10}
\]
Taking covariance of (10) with \(X_P\) and using (9) yields the normal
equations
\[
 (\Sigma_{c_0})_{Pj}
 =(\Sigma_{c_0})_{PP}B_{jP}^{\top}.                     \tag{11}
\]
Positive definiteness of \(\Sigma_{c_0}\) makes its principal submatrix
\((\Sigma_{c_0})_{PP}\) positive definite and therefore invertible.  Solving
(11) proves
\[
 B_{jP}^{\top}
 =(\Sigma_{c_0})_{PP}^{-1}(\Sigma_{c_0})_{Pj}.           \tag{12}
\]
Using (7), (9), and (10), or equivalently expanding the variance of the
residual in (10) and substituting (11)--(12), gives
\[
 \begin{aligned}
 \omega_j
 &=\operatorname{Var}(\varepsilon_j)\\
 &=(\Sigma_{c_0})_{jj}
   -(\Sigma_{c_0})_{jP}
    (\Sigma_{c_0})_{PP}^{-1}
    (\Sigma_{c_0})_{Pj}>0.                              \tag{13}
 \end{aligned}
\]
The strict inequality is also the positive Schur-complement property of the
positive-definite matrix \(\Sigma_{c_0}\).  If \(P=\varnothing\), the inverse
and products in (12)--(13) are interpreted as empty; then the coefficient row
has no parent entries and \(\omega_j=(\Sigma_{c_0})_{jj}>0\).

Finally consider the component matched to target \(j\).  If \(X^{(j)}\) has
its covariance and \(\eta^{(j)}=A^{(j)}X^{(j)}\), the perfect-surgery
factorization gives
\[
 \operatorname{Cov}(\eta^{(j)})
 =A^{(j)}\Sigma_j(A^{(j)})^{\top}
 =\Omega^{(j)}.                                         \tag{14}
\]
The \(j\)-th row of \(A^{(j)}\) is \(e_j^{\top}\), so
\(\eta_j^{(j)}=X_j^{(j)}\).  Taking the \((j,j)\) entry of (14) therefore
proves
\[
 \tau_j=(\Sigma_{\vartheta^{-1}(j)})_{jj}.              \tag{15}
\]

Equations (6), (12), (13), and (15) recover respectively \(G\), every
nonzero coefficient of \(B\), \(\omega\), and \(\tau\); zero entries of
\(B\) are fixed by the recovered parent sets.  The first step already
recovers the weight attached to each distinct covariance, so it also recovers
\(\pi\).  Thus the displayed operations define an observable recovery map on
the unordered covariance--weight collection, and the preceding equalities
prove, rather than define, that its value at every member of \(\mathcal M_p\)
is the generating
\(\theta_{\mathrm{sem}}=(c_0,\vartheta,G,B,\omega,\tau,\pi)\).  Composing
this map with the mixture identification in
\(\texttt{lem:component-distinctness}\) proves that \(P_{\mathrm{mix}}\)
uniquely determines \(\theta_{\mathrm{sem}}\), as claimed.
\end{proof}
\par\smallskip\noindent \textit{Justification.} The theorem identifies the baseline, targets, full DAG, coefficients, all variances, and attached mixture weights at every legal parameter point. \textit{Gap.} Squires et al. recover a broader latent representation under their setup but their unknown-baseline Proposition 10 is false. Kumar and Sinha (2021, Theorem 3.1) identify the intervention tuples and weights in a pooled mixture of perfect interventions only when the finite discrete baseline causal Bayesian network and its observational law are given, under positivity and exclusion. Kumar et al. provide the baseline separately for causal discovery and invoke interventional faithfulness. None identifies the hidden observational component and the full centered-Gaussian SEM explanation from the unordered mixture alone. \textit{Consumer.} Pooled perfect-intervention studies obtain a uniquely labeled structural Gaussian model without a designated control sample.

\begin{theorem}[thm:complete-certificate]\label{thm:complete-certificate}
The coordinate certificate \(\mathsf{Cert}\) accepts an arbitrary unordered \(p+1\)-tuple if and only if it has an explanation in \(\mathcal M_p\); on acceptance it returns the unique \(\theta_{\mathrm{sem}}\), and on rejection it returns a violated positive-definiteness, orientation, matching, baseline-factorization, or surgery equality. Its worst-case cost is \(O(p^4)\) exact real-arithmetic operations and \(O(p^3)\) storage.
\end{theorem}
\begin{proof}
We make explicit the checks named by \(\mathsf{Cert}\) in
\(\texttt{def:exact-certificate}\).  First check that the input consists
of \(p+1\) symmetric positive-definite, pairwise distinct matrices, and
that its weights are positive and sum to one.  Positive definiteness may
be certified or refuted by exact \(LDL^{\top}\) elimination, and it permits
all precision inversions below.

Scan the unordered precision pairs for one satisfying
\[
 K_c-K_d=s e_re_r^{\top},\qquad s\neq0.                            \tag{13}
\]
Use a third component \(z\) to designate as proposed baseline the endpoint
whose \(r\)-th covariance diagonal equals \((\Sigma_z)_{rr}\); reject at
the orientation check unless exactly one endpoint has this equality.
For the proposed baseline \(c_0\), form all original supports
\[
 S_c=\operatorname{rowsupp}(K_c-K_{c_0}).
\]
Maintain the unmatched-coordinate set \(U\).  At each step select a
remaining component for which \(S_c\cap U=\{j\}\), assign it to \(j\),
and remove \(c\) and \(j\).  Reject if some remaining residual support is
empty or if no singleton exists.  On completion put
\[
 P_j=S_{\vartheta^{-1}(j)}\setminus\{j\}.                           \tag{14}
\]
Because every member of \(P_j\) was removed before \(j\), the completed
trace orders every proposed edge \(i\to j\) forward and hence certifies
acyclicity.

Using the proposed baseline covariance, recover \(B_{jP_j}\) and
\(\omega_j\) by the regression and Schur-complement formulas (10)--(11),
and put
\[
 \tau_j=(\Sigma_{\vartheta^{-1}(j)})_{jj}.
\]
The Schur complement and every \(\tau_j\) are positive.  Finally check,
entry by entry,
\[
 K_{c_0}=A^{\top}\Omega^{-1}A                                    \tag{15}
\]
and, for every \(j\in V\),
\[
 K_{\vartheta^{-1}(j)}-K_{c_0}
 =\tau_j^{-1}e_je_j^{\top}-\omega_j^{-1}a_ja_j^{\top}.              \tag{16}
\]
The first failed preliminary, orientation, matching, factorization, or
surgery comparison supplies the advertised rejection trace; deterministic
lexicographic choices remove any implementation ambiguity.

\emph{Completeness.}
Suppose the tuple has an explanation in \(\mathcal M_p\).  The input and
weight checks pass by \(\texttt{def:model-class}\).
\(\texttt{lem:coordinate-pair-characterization}\) and the existence of a
root show that (13) is found, that any pair found has the baseline as an
endpoint, and that every third component orients it correctly.  By
\(\texttt{lem:surgery-identity}\), the supports are exactly (8).
At every stage a root of the DAG induced by the unmatched vertices has
singleton residual support, so matching never stalls or empties a support.
The uniqueness and regression argument in
\(\texttt{thm:universal-recovery}\) shows that the recovered target map,
parent sets, coefficients, and variances are the true ones.  Equations
(15)--(16) are then respectively the baseline factorization and the
surgery identity, so they pass.  Thus every legal tuple is accepted.

\emph{Soundness.}
Suppose the certificate accepts.  The matching trace makes
\(A=I_p-B\) the unit-diagonal equation matrix of an acyclic graph, hence
invertible.  Positivity of the recovered Schur complements makes
\(\Omega\) positive diagonal.  From (15),
\[
 \Sigma_{c_0}
 =(A^{\top}\Omega^{-1}A)^{-1}
 =A^{-1}\Omega A^{-\top},                                        \tag{17}
\]
so the proposed baseline has the required SEM factorization.

For the component matched to \(j\), replace row \(j\) of \(A\) by
\(e_j^{\top}\) and \(\omega_j\) by \(\tau_j\).  The resulting graph is a
subgraph of the accepted acyclic graph, so \(A^{(j)}\) remains invertible.
Expanding its precision row by row and applying (15)--(16) gives
\[
 \begin{aligned}
 (A^{(j)})^{\top}(\Omega^{(j)})^{-1}A^{(j)}
 &=A^{\top}\Omega^{-1}A
   -\omega_j^{-1}a_ja_j^{\top}
   +\tau_j^{-1}e_je_j^{\top}\\
 &=K_{\vartheta^{-1}(j)} .
 \end{aligned}                                                     \tag{18}
\]
Inverting (18) proves that every supplied nonbaseline covariance is
exactly the claimed perfect-surgery covariance.

It also follows from acceptance that the recovered graph and interventions
are effective.  If \(i\in P_j\) but \(B_{ji}=0\), the right side of (16)
has zero row \(i\), contradicting \(i\in S_{\vartheta^{-1}(j)}\).
Thus every asserted edge coefficient is nonzero.  If \(P_j=\varnothing\)
and \(\tau_j=\omega_j\), the right side of (16) is zero, contradicting the
checked distinctness of the component from the baseline.  Hence each
nonroot has a parent and every root changes variance.  Together with the
checked weights, (17)--(18) give an explanation in \(\mathcal M_p\).
By \(\texttt{thm:universal-recovery}\), this accepted explanation is
unique.

\emph{Cost.}
There are \(p+1\) dense \(p\times p\) matrices.  Exact
\(LDL^{\top}\) tests and inversions cost \(O(p^3)\) per matrix and hence
\(O(p^4)\) total.  There are \(O(p^2)\) unordered pairs, each of whose
\(p^2\) entries is inspected in (13), again costing \(O(p^4)\).
Computing baseline supports and naive singleton matching costs \(O(p^3)\).
Solving at most \(p\) dense principal regression systems costs at most
\(O(p^4)\).  Checking (15) costs \(O(p^3)\), and checking the \(p\)
outer-product identities (16) costs \(O(p^3)\).  Storing \(p+1\) dense
matrices costs \(O(p^3)\).  These are exact real-arithmetic operation and
storage bounds; no bit-complexity or finite-precision equality claim is
used.
\end{proof}
\par\smallskip\noindent \textit{Justification.} Identification alone cannot distinguish a legal perfect-surgery tuple from an algebraic impostor; soundness and completeness make the reconstruction semantically auditable. \textit{Gap.} Appendix G of Squires et al. explicitly leaves full membership testing open and tests only a necessary rank condition. Kumar and Sinha (2021, Theorem 3.1) recover intervention tuples relative to a given finite discrete baseline causal Bayesian network and observational law; they do not certify whether an arbitrary unordered Gaussian component tuple admits one common hidden-baseline SEM and all claimed perfect surgeries. \textit{Consumer.} A pooled-data pipeline can reject component collections that cannot arise from one common baseline DAG before invoking causal interpretation.

\begin{theorem}[thm:rank-selector-correction]\label{thm:rank-selector-correction}
Theorem 2 strictly replaces the unknown-baseline promise of Proposition 10 in Squires et al. on the observed-coordinate fully covered perfect-intervention subclass: both two-node constructions are legal members of \(\mathcal M_2\), the first gives literal scores \((3,2,2)\) and symmetric scores \((3,3,4)\), the second gives symmetric scores \((3,2,3)\), and the legal three-node fork gives literal scores \((5,4,5,5)\); in each case the coordinate certificate returns the true baseline and explanation.
\end{theorem}
\begin{proof}
By \(\texttt{lem:rank-witness-calculations}\), both displayed two-node
chains satisfy all defining conditions of \(\mathcal M_2\).  The first
has literal scores \((3,2,2)\) and symmetric scores \((3,3,4)\); the
second has symmetric scores \((3,2,3)\).  The same lemma proves that the
three-node fork is legal and has literal scores \((5,4,5,5)\).  Hence the
literal rule excludes the true baseline in the first chain and the fork,
the symmetric rule ties the baseline with another context in the first
chain, and the symmetric rule strictly prefers the root intervention in
the second chain.

The full-domain force of these calculations follows from
\(\texttt{prop:squires-rank-failure-full-domain}\), which in turn uses the
verified reverse-topological embedding of
\(\texttt{lem:squires-class-embedding}\).  Thus the negative conclusion
is not an artifact of imposing identity latent labels or of leaving the
published model class.

For the positive statement, each displayed covariance tuple is a member
of \(\mathcal M_p\).  The completeness and uniqueness clauses of
\(\texttt{thm:complete-certificate}\) therefore imply that
\(\mathsf{Cert}\) accepts it and returns its true baseline and full
semantic explanation.  This establishes the stated strict replacement:
the rank selector is invalid on the full published domain, whereas the
coordinate certificate is valid on the observed-coordinate, fully covered
perfect-surgery subclass \(\mathcal M_p\).  No positive recovery claim is
made for the rest of the latent-mixing class.
\end{proof}
\par\smallskip\noindent \textit{Justification.} The reverse-topologically embedded low-dimensional and open star families locate the failed rank inequality on the full published domain, while the coordinate certificate supplies the valid positive replacement on the observed-coordinate fully covered subclass. \textit{Gap.} Proposition 10's unique rank-score minimization fails inside its full published class; the repaired permutation embedding establishes that transfer without asserting positive recovery under latent mixing. \textit{Consumer.} Users of the ICML reduction obtain exact regression tests for the failed rank selector and a valid baseline preprocessing certificate when their covariance tuple satisfies the observed-coordinate fully covered model.

\begin{theorem}[thm:rank-only-impossibility]\label{thm:rank-only-impossibility}
For every \(p\ge2\), the one-root star construction contains a nonempty strict-inequality parameter region on which the pairwise rank matrix \(R\) is invariant under swapping the baseline with the root-intervention vertex. Consequently, no component-permutation-equivariant rank-only selector universally recovers the baseline on \(\mathcal M_p\).
\end{theorem}
\begin{proof}
Fix \(p\geq 2\) and a member of the one-root star family in \(\texttt{def:star-counterfamily}\).  Thus vertex \(1\) is the unique root and, for each \(j\in\{2,\ldots,p\}\), the only edge into \(j\) is \(1\to j\), with coefficient \(b_j\neq0\).  From \(A=I_p-B\) and \(a_j=A^{\top}e_j\), the transposes of the structural-equation rows are
\[
 a_1=e_1,
 \qquad
 a_j=e_j-b_je_1\quad (j=2,\ldots,p).
 \tag{1}
\]
All \(\omega_j\) and \(\tau_j\) are strictly positive.  Define
\[
 d=\tau_1^{-1}-\omega_1^{-1}.
\]
The star-family restriction gives \(d\neq0\); this is also the effectiveness condition for the root intervention because \(\operatorname{Pa}(1)=\varnothing\).  Applying \(\texttt{lem:surgery-identity}\) to the root and using (1) gives
\[
 \Delta_1=d e_1e_1^{\top}.
 \tag{2}
\]
For a leaf \(j\geq2\), the same lemma and (1) give
\[
 \Delta_j
 =\tau_j^{-1}e_je_j^{\top}
  -\omega_j^{-1}(e_j-b_je_1)(e_j-b_je_1)^{\top}.
 \tag{3}
\]
Therefore \(\Delta_j\) vanishes outside rows and columns \(\{1,j\}\), and its principal submatrix on those coordinates is
\[
 \left.\Delta_j\right|_{\{1,j\}}
 =
 \begin{pmatrix}
 -b_j^2/\omega_j & b_j/\omega_j\\
 b_j/\omega_j & \tau_j^{-1}-\omega_j^{-1}
 \end{pmatrix}.
 \tag{4}
\]
Direct expansion, using positivity of \(\omega_j\) and \(\tau_j\), yields
\[
 \begin{aligned}
 \det\!\left(\left.\Delta_j\right|_{\{1,j\}}\right)
 &= -\frac{b_j^2}{\omega_j}
       \left(\frac1{\tau_j}-\frac1{\omega_j}\right)
    -\frac{b_j^2}{\omega_j^2}\\
 &= -\frac{b_j^2}{\omega_j\tau_j}<0,
 \end{aligned}
 \tag{5}
\]
where strictness also uses \(b_j\neq0\).  Since all other rows and columns vanish, (5) proves
\[
 R_{0j}=\operatorname{rank}(K_j-K_0)
 =\operatorname{rank}(\Delta_j)=2.
 \tag{6}
\]
Subtracting (2) from (4) and expanding again gives
\[
 \begin{aligned}
 \det\!\left(
   \left.(\Delta_j-\Delta_1)\right|_{\{1,j\}}
 \right)
 &=\left(-\frac{b_j^2}{\omega_j}-d\right)
      \left(\frac1{\tau_j}-\frac1{\omega_j}\right)
      -\frac{b_j^2}{\omega_j^2}\\
 &=\frac{d(\tau_j-\omega_j)-b_j^2}{\omega_j\tau_j}.
 \end{aligned}
 \tag{7}
\]
The defining strict inequality in \(\texttt{def:star-counterfamily}\) makes (7) nonzero.  Hence
\[
 R_{1j}=\operatorname{rank}(K_j-K_1)
 =\operatorname{rank}(\Delta_j-\Delta_1)=2.
 \tag{8}
\]
Finally, (2) and \(d\neq0\) imply
\[
 R_{01}=\operatorname{rank}(K_1-K_0)=1.
 \tag{9}
\]

Let \(\sigma\) denote only the transposition of component labels \(0\) and \(1\), fixing every leaf-component label, and let \(P_\sigma\) be its permutation matrix on \(C\).  Equations (6) and (8) give \(R_{0j}=R_{1j}\) for every \(j\geq2\); equation (9) is symmetric in \(0\) and \(1\); entries indexed by two leaves are unchanged because \(\sigma\) fixes both indices; and every diagonal entry is zero.  These cases exhaust \(C\times C\), so
\[
 P_\sigma R P_\sigma^{\top}=R.
 \tag{10}
\]

The strict region is nonempty.  Set
\[
 \omega_k=1\quad(k\in V),\qquad
 \tau_1=2,\qquad
 \tau_j=1\quad(j=2,\ldots,p),\qquad
 b_j=1\quad(j=2,\ldots,p),
 \tag{11}
\]
and choose, for example, \(\pi_c=1/(p+1)\) for every \(c\in C\).  Then all variances and weights are positive, the weights sum to one, \(d=-1/2\neq0\), and for each leaf
\[
 d(\tau_j-\omega_j)-b_j^2=-1<0.
 \tag{12}
\]
The root intervention is effective because \(\tau_1\neq\omega_1\); every leaf intervention is effective because it has the parent \(1\), even though (11) has \(\tau_j=\omega_j\).  Together with the acyclic star graph and the baseline and surgery equations, these facts place the example in \(\mathcal M_p\) by \(\texttt{def:model-class}\).  Positivity, \(d\neq0\), the conditions \(b_j\neq0\), and the finitely many strict inequalities in (12) persist on a sufficiently small Euclidean neighborhood of the displayed variance and coefficient parameters; positivity of the weights persists on a relative neighborhood in the probability simplex.  Hence the family contains a nonempty relatively open strict-inequality parameter region.

For completeness, all \(p+1\) Gaussian components are distinct on this region.  Equation (2) separates the root intervention from the baseline; (5) separates every leaf intervention from the baseline; (7) separates every leaf intervention from the root intervention.  If \(j\neq k\) are leaves, (3) gives \((\Delta_j)_{1j}=b_j/\omega_j\neq0\), whereas \((\Delta_k)_{1j}=0\), so the two leaf components are distinct as well.

Now let \(F\) be any component-permutation-equivariant rank-only baseline selector.  By \(\texttt{def:rank-compression}\), it is a map from rank matrices to component labels satisfying
\[
 F(P_\gamma R P_\gamma^{\top})=\gamma(F(R))
 \tag{13}
\]
for every permutation \(\gamma\) of \(C\).  Apply (13) with \(\gamma=\sigma\) and use (10):
\[
 F(R)=\sigma(F(R)).
 \tag{14}
\]
Neither \(0\) nor \(1\) is fixed by \(\sigma\).  Thus (14) forces \(F(R)\in C\setminus\{0,1\}\); this remains valid for \(p=2\), when the fixed-point set is the singleton \(\{2\}\).  In particular, \(F(R)\neq0\), while component \(0\) is the observational baseline.  Therefore \(F\) fails at every point of the nonempty open region.  Since \(F\) was arbitrary, no component-permutation-equivariant rank-only selector universally recovers the baseline on \(\mathcal M_p\).

The same obstruction also transfers, negatively, to the full model class of Squires et al., without any positive recovery claim under latent mixing.  Indeed, \(\texttt{lem:squires-class-embedding}\) constructs a distinct reverse-topological latent-node bijection \(\rho\) and permutation matrix \(Qe_j=e_{\rho(j)}\), verifies the published assumptions, leaves the context labels unchanged, and proves \(\Theta_c=K_c\) for every \(c\in C\).  It therefore preserves every entry of \(R\).  The symbol \(\sigma\) in (10)--(14) remains solely the component-label transposition and is not the latent relabeling \(\rho\).  Consequently, any equivariant rank-only selector universally correct on the full Proposition-10 class would restrict to a selector correct on this embedded open star family, contradicting (14).
\end{proof}
\par\smallskip\noindent \textit{Justification.} The obstruction shows that the correction needs entrywise support and marginal information, rather than a different scalar aggregation of the same ranks. \textit{Gap.} Squires et al. use one rank score, but the literature does not establish an open-family impossibility for the entire class of equivariant rank-only selectors. \textit{Consumer.} Method designers can rule out every rank-only baseline heuristic and focus on information-preserving summaries.

\begin{theorem}[thm:stable-discrete-recovery]\label{thm:stable-discrete-recovery}
If \(E_{\mathrm{op}}<\varepsilon_{\star}\), then after permutation matching every \(\widehat\Sigma_c\) is positive definite, every precision-difference entry has error below \(\alpha/4\), and the thresholded coordinate decoder recovers \(c_0\), \(\vartheta\), and \(G\) exactly. The condition uses the minimum eigenvalue across all components, not only the baseline.
\end{theorem}
\begin{proof}
Because the permutation group is finite, the minimum in \(E_{\mathrm{op}}\)
is attained.  Relabel the estimates by an attaining permutation and put
\[
 \varepsilon=\max_{c\in C}
 \|\widehat\Sigma_c-\Sigma_c\|_{\mathrm{op}}
 <\varepsilon_\star .
\]
Every true covariance has minimum eigenvalue at least \(\lambda\).  For
every unit vector \(x\), the operator-norm bound gives
\[
 x^{\top}\widehat\Sigma_cx
 =x^{\top}\Sigma_cx+
   x^{\top}(\widehat\Sigma_c-\Sigma_c)x
 \ge\lambda-\varepsilon.
\]
Taking the minimum over unit vectors and using
\(\varepsilon<\lambda/2\) gives
\[
 \lambda_{\min}(\widehat\Sigma_c)
 \ge\lambda_{\min}(\Sigma_c)-\varepsilon
 \ge\lambda-\varepsilon>\lambda/2>0.                               \tag{25}
\]
Thus every estimated covariance is positive definite and invertible.
The resolvent identity gives
\[
 \widehat K_c-K_c
 =\widehat K_c(\Sigma_c-\widehat\Sigma_c)K_c.
\]
Using (25),
\[
 \|\widehat K_c-K_c\|_{\mathrm{op}}
 \le\frac{\varepsilon}{\lambda(\lambda-\varepsilon)}
 <\frac{2\varepsilon}{\lambda^2}.                                 \tag{26}
\]
Consequently, for every pair \(c,d\) and every entry,
\[
 \left|
 [(\widehat K_c-\widehat K_d)-(K_c-K_d)]_{ik}
 \right|
 \le\|\widehat K_c-K_c\|_{\mathrm{op}}
    +\|\widehat K_d-K_d\|_{\mathrm{op}}
 <\frac{4\varepsilon}{\lambda^2}
 <\frac{\alpha}{4}.                                                \tag{27}
\]
This proves the stated precision-difference error bound.  In particular,
a true zero has estimated magnitude below \(\alpha/4\), while every
nonzero entry of any \(\Delta_j\) has magnitude at least \(\alpha\) and
is estimated with magnitude greater than \(3\alpha/4\).  Thresholding at
\(\alpha/2\) therefore preserves every zero/nonzero decision needed below.

We next prove that the detected coordinate pairs are exactly the
baseline--root pairs.  Such a true pair has one nonzero diagonal entry
and no other entry, so (27) makes it detected exactly.  A
baseline--nonroot pair has a nonzero parent--target off-diagonal entry of
magnitude at least \(\alpha\), so it is rejected.  For two intervention
targets \(u,v\), order them so that \(u\) precedes \(v\) topologically.
Row \(v\) of \(\Delta_u\) is zero.  If \(v\) is nonroot, a nonzero
parent--target off-diagonal entry in row \(v\) of \(\Delta_v\) survives in
\(\Delta_v-\Delta_u\), so the pair is rejected.  If \(v\) is a root and
\(u\) is nonroot, a nonzero off-diagonal entry of \(\Delta_u\) survives
the diagonal matrix \(\Delta_v\).  If both are roots, their difference
has two nonzero diagonal entries on distinct coordinates.  By (27), the
surviving witnesses are detected in every case.  Hence no false pair is
recognized.

Consider a detected root pair at coordinate \(r\), and any third
component \(z\).  At truth,
\[
 (\Sigma_0)_{rr}=(\Sigma_z)_{rr}=\omega_r,
 \qquad
 (\Sigma_r)_{rr}=\tau_r,
 \qquad |\tau_r-\omega_r|\ge g .
\]
An operator-norm covariance error bounds each diagonal-entry error.
Therefore
\[
 |(\widehat\Sigma_0)_{rr}-(\widehat\Sigma_z)_{rr}|
 \le2\varepsilon<g/4,
\]
whereas
\[
 |(\widehat\Sigma_r)_{rr}-(\widehat\Sigma_z)_{rr}|
 \ge g-2\varepsilon>3g/4.                                         \tag{28}
\]
The closer endpoint in \(\texttt{def:threshold-decoder}\) is consequently
the true baseline, for every permitted third component.

After this orientation, apply the same \(\alpha/2\) threshold to
\(\widehat K_c-\widehat K_{c_0}\).  By (27), every zero row of
\(\Delta_j\) remains undetected, and every nonzero row contains an entry
of magnitude at least \(\alpha\) and is detected.  Hence all thresholded
row supports equal
\[
 \{j\}\cup\operatorname{Pa}(j)
\]
by \(\texttt{lem:surgery-identity}\).  The unique matching argument in
\(\texttt{thm:universal-recovery}\) then recovers \(\vartheta\), and
removing each matched target from its support recovers every parent set
and therefore \(G\).  Undoing the analytical relabeling maps \(c_0\) and
\(\vartheta\) back to the original estimated-component positions.  The
decoder itself does not need to know the attaining permutation.

Finally, all margins used above are positive: \(\lambda>0\) because the
finite collection consists of positive-definite covariances;
\(\alpha>0\) because finitely many nonzero entries are minimized and every
\(\Delta_j\neq0\); and \(g>0\) because every finite DAG has a root and
every root intervention changes variance.  Thus every division and
threshold comparison in the proof is valid.
\end{proof}
\par\smallskip\noindent \textit{Justification.} The resolvent radius converts population identification into a deterministic error interface for any component estimator. \textit{Gap.} The cited causal-discovery papers do not give an all-component operator-norm neighborhood guaranteeing unknown-baseline and exact target recovery without faithfulness. \textit{Consumer.} Any future mixture estimator with a permutation-matched operator tail inherits a finite-sample discrete-recovery bound by event inclusion.

\begin{proposition}[prop:baseline-only-failure]\label{prop:baseline-only-failure}
In the near-singular family, the baseline minimum eigenvalue is fixed, \(\alpha=1/2\), and \(g>1/2\), yet \(\|\widehat\Sigma_1(t)-\Sigma_1(t)\|_{\mathrm{op}}=t\to0\) and every estimated precision pair has a detected off-diagonal entry. At \(t=1/100\), the three off-diagonal magnitudes are \(103/97\), \(1\), and \(200/97\). Thus no positive radius depending only on baseline conditioning, \(\alpha\), and \(g\) can certify this inverse-threshold decoder.
\end{proposition}
\begin{proof}
Fix \(0<t<1/3\).  We first verify that the displayed covariance triple is a legal interior point of the stated causal model, rather than merely a collection of positive-definite matrices.  On coordinates \(1,2\), take the acyclic graph \(1\mathbin{\to}2\) and the baseline structural equations
\[
 X_1=\varepsilon_1,\qquad X_2=X_1+\varepsilon_2,
 \qquad \varepsilon_1,\varepsilon_2\ \text{independent with law }N(0,1).
\]
They give
\[
 \operatorname{Cov}(X_1,X_2)
 =\begin{pmatrix}1&1\\1&2\end{pmatrix}=\Sigma_0.
\]
Replacing only the first equation by an independent draw \(X_1\sim N(0,t)\), while retaining the second equation, gives \(\Sigma_1(t)\).  Replacing only the second equation by an independent draw \(X_2\sim N(0,2)\), while retaining the first equation, gives \(\Sigma_2\).  Thus the baseline noise variances and intervention variances are respectively
\[
 \omega=(1,1),\qquad \tau=(t,2).
\]
Both interventions are effective because \(t\ne1\) and \(2\ne1\).  The unique root is coordinate \(1\), so Definition~\(\mathrm{def:stability\text{-}margins}\) gives
\[
 g=|\tau_1-\omega_1|=1-t>\frac23>\frac12.
\]
All three components are distinct.  Moreover, their first diagonal entries and determinants are
\[
 \begin{array}{c|ccc}
  &\Sigma_0&\Sigma_1(t)&\Sigma_2\\ \hline
 \text{first diagonal entry}&1&t&1\\
 \det&1&t&2,
 \end{array}
\]
For a symmetric matrix \(Q=\left(\begin{smallmatrix}q_{11}&q_{12}\\q_{12}&q_{22}\end{smallmatrix}\right)\) with \(q_{11}>0\), direct completion of the square gives
\[
 (x,y)Q(x,y)^{\mathsf T}
 =q_{11}\left(x+\frac{q_{12}}{q_{11}}y\right)^2
  +\frac{\det Q}{q_{11}}y^2.
\]
The displayed positive entries and determinants therefore prove directly that each covariance is positive definite.  Consequently every component is a nonsingular centered Gaussian with support \(\mathbb R^2\); arbitrary fixed strictly positive mixing weights summing to one may be attached without affecting the calculation.

The characteristic polynomial of \(\Sigma_0\) is \(x^2-3x+1\).  Hence its minimum eigenvalue is the fixed positive number
\[
 \lambda_{0}=\lambda_{\min}(\Sigma_0)=\frac{3-\sqrt5}{2},
\]
independent of \(t\).  Direct inversion yields
\[
 K_0=\begin{pmatrix}2&-1\\-1&1\end{pmatrix},\qquad
 K_1(t)=\begin{pmatrix}t^{-1}+1&-1\\-1&1\end{pmatrix},\qquad
 K_2=\begin{pmatrix}1&0\\0&1/2\end{pmatrix}.
\]
Therefore, using the precision-difference convention in Definition~\(\mathrm{def:stability\text{-}margins}\),
\[
 \Delta_1=K_1-K_0=
 \begin{pmatrix}t^{-1}-1&0\\0&0\end{pmatrix},
 \qquad
 \Delta_2=K_2-K_0=
 \begin{pmatrix}-1&1\\1&-1/2\end{pmatrix}.
\]
Since \(t^{-1}-1>2\), the minimum absolute value among all nonzero entries of \(\Delta_1\) and \(\Delta_2\) is exactly
\[
 \alpha=\frac12.
\]

Now leave \(\Sigma_0\) and \(\Sigma_2\) unchanged and use the perturbation in Definition~\(\mathrm{def:near\text{-}singular\text{-}family}\):
\[
 E_t:=\widehat\Sigma_1(t)-\Sigma_1(t)
 =\begin{pmatrix}0&t\\t&0\end{pmatrix}.
\]
The eigenvalues of \(E_t\) are \(t\) and \(-t\), whence
\[
 \|\widehat\Sigma_1(t)-\Sigma_1(t)\|_{\mathrm{op}}=t\longrightarrow0
 \quad\text{as }t\downarrow0.
\]
This limit is taken through \(t>0\); the singular boundary point \(t=0\) is never included.  The perturbed matrix is itself positive definite by the same completed-square identity because its first diagonal entry is \(t>0\) and
\[
 \det\widehat\Sigma_1(t)=t(1+t)-4t^2=t(1-3t)>0.
\]
Its inverse is
\[
 \widehat K_1(t)
 =\frac1{t(1-3t)}
 \begin{pmatrix}1+t&-2t\\-2t&t\end{pmatrix}.
\]
In particular, \((\widehat K_1)_{12}=-2/(1-3t)\), whereas \((K_0)_{12}=-1\) and \((K_2)_{12}=0\).  Thus the absolute off-diagonal entries in the three estimated pairwise precision differences are, respectively,
\[
 \left| (\widehat K_1-K_0)_{12}\right|
   =\frac{1+3t}{1-3t},\qquad
 \left| (K_2-K_0)_{12}\right|=1,
 \qquad
 \left| (\widehat K_1-K_2)_{12}\right|
   =\frac2{1-3t}.
\]
Every denominator is positive, and these three quantities are all strictly greater than \(1/4=\alpha/2\).  By Definition~\(\mathrm{def:threshold\text{-}decoder}\), each pair therefore has a detected off-diagonal entry.  A coordinate pair must have one detected diagonal entry and no other detected entry, so none of the three pairs is a coordinate pair.  The decoder consequently cannot recover its required root pair, regardless of the input ordering.

For the asserted uniform impossibility, interpret the advertised margins in their usual certificate sense as lower bounds.  Suppose a baseline-only theorem assigned a positive uniform radius \(r\) to the fixed baseline conditioning \(\lambda_{\min}(\Sigma_0)=\lambda_0\), the support margin lower bound \(\alpha\ge1/2\), and the root-gap lower bound \(g\ge1/2\).  Choose
\[
 0<t<\min\{1/3,r\}.
\]
The preceding legal model has exactly the fixed baseline covariance, has \(\alpha=1/2\) and \(g>1/2\), and satisfies
\[
 \max_{c\in\{0,1,2\}}
 \|\widehat\Sigma_c-\Sigma_c\|_{\mathrm{op}}=t<r,
\]
yet its threshold decoder fails.  Since this construction works for every \(r>0\), no positive uniform radius based only on baseline conditioning and those fixed structural margins can certify this inverse-threshold decoder.  The obstruction is precisely the unmonitored limit \(\lambda_{\min}(\Sigma_1(t))\le e_1^{\mathsf T}\Sigma_1(t)e_1=t\to0\).

Finally, at \(t=1/100\), substitution into the three displayed off-diagonal magnitudes gives
\[
 \frac{1+3/100}{1-3/100}=\frac{103}{97},\qquad
 1,\qquad
 \frac{2}{1-3/100}=\frac{200}{97},
\]
as claimed.
\end{proof}
\par\smallskip\noindent \textit{Justification.} This boundary result explains why the stability theorem conditions every component rather than concealing near-singular interventions. \textit{Gap.} No cited result isolates this conditioning failure for unknown-baseline perfect-intervention decoding. \textit{Consumer.} Finite-sample implementations learn which eigenvalue margin must be monitored before reporting a stable causal explanation.

\begin{proposition}[prop:factorization-necessity]\label{prop:factorization-necessity}
The three-dimensional false-positive tuple passes singleton support matching and every local diagonal surgery check for an apparent empty DAG, but fails \(K_{c_0}=A^{\top}\Omega^{-1}A\); it therefore has no explanation in \(\mathcal M_3\).
\end{proposition}
\begin{proof}
Put \(E_j=e_je_j^{\top}\).  By \(\mathrm{def:false\mbox{-}positive\mbox{-}tuple}\),
\[
 K_0=\begin{pmatrix}1&0&0\\0&2&-1\\0&-1&2\end{pmatrix},
 \qquad K_j=K_0+E_j\quad(j\in\{1,2,3\}).
\]
These are legitimate precision matrices.  Indeed, for every nonzero
\(x=(x_1,x_2,x_3)^{\top}\),
\[
 x^{\top}K_0x=x_1^2+x_2^2+x_3^2+(x_2-x_3)^2>0,
\]
and \(x^{\top}K_jx=x^{\top}K_0x+x_j^2>0\).  Direct inversion gives
\[
 \begin{aligned}
 \Sigma_0&=\begin{pmatrix}1&0&0\\0&2/3&1/3\\0&1/3&2/3\end{pmatrix},&
 \Sigma_1&=\begin{pmatrix}1/2&0&0\\0&2/3&1/3\\0&1/3&2/3\end{pmatrix},\\
 \Sigma_2&=\begin{pmatrix}1&0&0\\0&2/5&1/5\\0&1/5&3/5\end{pmatrix},&
 \Sigma_3&=\begin{pmatrix}1&0&0\\0&3/5&1/5\\0&1/5&2/5\end{pmatrix}.
 \end{aligned}
\]

First run the local parts of \(\mathsf{Cert}\) from
\(\mathrm{def:exact\mbox{-}certificate}\) with proposed baseline \(c_0=0\).  For
each \(j\in V\),
\[
 \Delta_j:=K_j-K_0=E_j,
 \qquad \operatorname{rowsupp}(\Delta_j)=\{j\}.
\]
Thus singleton matching assigns component \(j\) to target \(j\) and returns
\(\operatorname{Pa}(j)=\varnothing\) for every \(j\), hence \(A=I_3\).
Baseline regression with no regressors has residual variance equal to the
corresponding marginal variance.  The reconstructed variances are therefore
\[
 \omega=(1,2/3,2/3),
 \qquad
 \tau=((\Sigma_1)_{11},(\Sigma_2)_{22},(\Sigma_3)_{33})=(1/2,2/5,2/5).
\]
All three apparent root interventions are effective because
\(\tau_j\neq\omega_j\).  Since the equation-row vectors are
\(a_j=e_j\), the three local surgery identities required by the certificate
hold entrywise:
\[
 \tau_j^{-1}E_j-\omega_j^{-1}a_ja_j^{\top}
 =\bigl(\tau_j^{-1}-\omega_j^{-1}\bigr)E_j
 =E_j
 =K_j-K_0,
 \qquad j=1,2,3,
\]
where the three scalar differences are
\(2-1=1\), \(5/2-3/2=1\), and \(5/2-3/2=1\).
Nevertheless, the common-baseline factorization fails, since
\[
 K_0-A^{\top}\Omega^{-1}A
 =K_0-\operatorname{diag}(1,3/2,3/2)
 =\begin{pmatrix}0&0&0\\0&1/2&-1\\0&-1&1/2\end{pmatrix}\neq0.
\]
This proves all of the asserted local passes and the asserted failure of the
global check.

It remains to show that choosing another member of the unordered collection
as baseline cannot yield a legal explanation.  Suppose, to the contrary, that
the collection has an explanation in \(\mathcal M_3\), as defined by
\(\mathrm{def:model\mbox{-}class}\), and let \(b\in\{0,1,2,3\}\) denote the
component used as its baseline.  Acyclicity makes the unit-diagonal equation
matrix invertible: after a topological permutation it is triangular.  Inverting
the baseline covariance factorization and each intervened covariance
factorization therefore yields
\[
 K_b=A^{\top}\Omega^{-1}A,
 \qquad
 K_c-K_b=\tau_t^{-1}E_t-\omega_t^{-1}a_ta_t^{\top}
\]
when the nonbaseline component \(c\) is assigned to intervention target
\(t\).  In equation-row convention,
\[
 a_t=e_t-\sum_{i\in\operatorname{Pa}(t)}B_{ti}e_i.
\]
Consequently no row outside
\(\{t\}\cup\operatorname{Pa}(t)\) can be nonzero in the displayed precision
difference.  If \(i\in\operatorname{Pa}(t)\), its \((i,t)\) entry equals
\(B_{ti}/\omega_t\neq0\).  If \(t\) is a root, the difference equals
\((\tau_t^{-1}-\omega_t^{-1})E_t\), which is nonzero by the effective-root
condition in \(\mathrm{def:model\mbox{-}class}\).  Hence, in all cases,
\[
 \operatorname{rowsupp}(K_c-K_b)=\{t\}\cup\operatorname{Pa}(t).
 \tag{1}
\]

If \(b=0\), the three observed supports are the singletons
\(\{1\},\{2\},\{3\}\).  Equation \((1)\), together with the bijection between
the three nonbaseline components and the three targets, forces component
\(j\) to target \(j\) and forces every parent set to be empty.  Thus
\(A=I_3\), but then \(K_b=A^{\top}\Omega^{-1}A\) is diagonal, contrary to
\((K_0)_{23}=-1\).

Now let \(b=r\in\{1,2,3\}\).  Relative to this proposed baseline the three
nonbaseline differences are
\[
 K_0-K_r=-E_r,
 \qquad K_s-K_r=E_s-E_r\quad(s\in V\setminus\{r\}),
\]
with row supports \(\{r\}\) and \(\{r,s\}\), respectively.  Equation \((1)\)
forces component \(0\) to target \(r\).  Bijectivity then forces component
\(s\) to target \(s\) for each \(s\neq r\), and it forces
\[
 \operatorname{Pa}(r)=\varnothing,
 \qquad \operatorname{Pa}(s)=\{r\}\quad(s\neq r).
 \tag{2}
\]
Thus the only possible graph is the directed star rooted at \(r\), with both
edge coefficients nonzero.

For \(r=1\), the three equation rows have supports contained respectively in
\(\{1\}\), \(\{1,2\}\), and \(\{1,3\}\).  Therefore the factorization
\(K_1=\sum_{q=1}^3\omega_q^{-1}a_qa_q^{\top}\) implies
\((K_1)_{23}=0\), whereas the displayed tuple has \((K_1)_{23}=-1\).
For \(r=2\), condition \((2)\) gives \(B_{12}\neq0\) and
\[
 (K_2)_{12}=-B_{12}/\omega_1\neq0,
\]
because no equation row other than
\(a_1=e_1-B_{12}e_2\) contributes to that entry; but the displayed tuple has
\((K_2)_{12}=0\).  Finally, for \(r=3\), condition \((2)\) gives
\(B_{13}\neq0\) and analogously
\[
 (K_3)_{13}=-B_{13}/\omega_1\neq0,
\]
whereas the displayed tuple has \((K_3)_{13}=0\).  Every possible baseline
component leads to a contradiction.  Hence the tuple has no explanation in
\(\mathcal M_3\).  Arbitrary positive component weights summing to one do not
alter any step and satisfy the weight requirement in
\(\mathrm{def:model\mbox{-}class}\).
\end{proof}
\par\smallskip\noindent \textit{Justification.} The example makes the baseline factorization check logically necessary for certificate soundness. \textit{Gap.} The necessary rank screen in Squires et al. does not reject all tuples whose local differences mimic atomic interventions. \textit{Consumer.} Certificate implementers receive a minimal false-positive test that catches omission of the global common-baseline constraint.

\begin{openendedquestion}[oeq:growing-mixture-tail]\label{oeq:growing-mixture-tail}
Using the growing-component estimation handle, construct a concrete estimator from \(Z_1,\ldots,Z_n\) and derive an explicit \(T_n(\delta)\) such that \(\Pr\{E_{\mathrm{op}}>T_n(\delta)\}\le\delta\) when the number of centered Gaussian components is \(p+1\); characterize primitive weight, spectral-envelope, and pairwise-separation conditions under which \(T_n(\delta)<\varepsilon_{\star}\) with sample size polynomial in \(p\), and determine whether the perfect-surgery geometry weakens the near-disjoint total-variation separation required by generic centered-mixture clustering.
\end{openendedquestion}
\par\smallskip\noindent \textit{Justification.} This is the unresolved end-to-end statistical bottleneck rather than an assumed component-estimation oracle. \textit{Gap.} Kumar et al. fix the component count, while Anderson et al. cluster growing centered mixtures under strong total-variation separation but do not deliver the surgery-specific operator-error tail consumed here. \textit{Consumer.} A solution would turn the deterministic certificate into a growing-dimensional guarantee for genuinely pooled intervention samples.

\begin{lemma}[lem:squires-class-embedding]\label{lem:squires-class-embedding}
Every member of \(\mathcal M_p\) admits an embedding into the full model class of Squires et al., Appendix F.2, Proposition 10. Choose a reverse-topological bijection \(\rho:V\to V\) of the baseline DAG, let \(Qe_j=e_{\rho(j)}\) with \(Q^{\top}Q=I_p\), and put \(\mathcal C_0=\Omega^{-1/2}A\), \(\mathcal C_j=(\Omega^{(j)})^{-1/2}A^{(j)}=\mathcal C_0+e_j(\tau_j^{-1/2}e_j-\mathcal C_0^{\top}e_j)^{\top}\), \(d=p\), \(G_{\mathrm{Sq}}=Q^{\top}\), \(H_{\mathrm{Sq}}=G_{\mathrm{Sq}}^{\dagger}=Q\), and \(\mathcal B_c^{\mathrm{Sq}}=Q\mathcal C_cQ^{\top}\) for \(c\in C\). Then context \(j\) targets latent node \(\rho(j)\), the reverse-topological, generic-row, and row-normalization requirements of the published class hold, and \[ \mathcal B_j^{\mathrm{Sq}}=\mathcal B_0^{\mathrm{Sq}}+e_{\rho(j)}\bigl(\tau_j^{-1/2}e_{\rho(j)}-(\mathcal B_0^{\mathrm{Sq}})^{\top}e_{\rho(j)}\bigr)^{\top}. \] The context labels themselves are unchanged, and the published context precisions satisfy \[ \Theta_c=H_{\mathrm{Sq}}^{\top}(\mathcal B_c^{\mathrm{Sq}})^{\top}\mathcal B_c^{\mathrm{Sq}}H_{\mathrm{Sq}}=K_c \qquad(c\in C). \] Thus the embedding preserves every pairwise precision-difference rank, the literal Proposition-10 score, and the symmetric score.
\end{lemma}
\begin{proof}
Fix an arbitrary member of \(\mathcal M_p\).  All diagonal entries of
\(\Omega\) and \(\Omega^{(j)}\) are positive by
\(\texttt{def:model-class}\), so their inverse square roots exist.  A
finite DAG has a reverse-topological bijection: inductively, choose a root,
assign it the largest unused label, delete it, and order the remaining DAG.
Thus there is a bijection \(\rho:V\to V\) such that
\[
 u\to v\quad\Longrightarrow\quad \rho(u)>\rho(v).                 \tag{1}
\]
Let \(Q\) be the permutation matrix determined by
\(Qe_j=e_{\rho(j)}\).  Then
\(Q^{-1}=Q^{\top}\) and \(Q^{\top}Q=I_p\).

Define the normalized equation matrices
\[
 \mathcal C_0=\Omega^{-1/2}A,
 \qquad
 \mathcal C_j=(\Omega^{(j)})^{-1/2}A^{(j)}
 \quad(j\in V).                                                  \tag{2}
\]
For a fixed \(j\), every row of the second matrix in (2), other than row
\(j\), equals the corresponding row of \(\mathcal C_0\), while row
\(j\) is \(\tau_j^{-1/2}e_j^{\top}\).  Hence, with
\[
 c_j=\tau_j^{-1/2}e_j-\mathcal C_0^{\top}e_j,
\]
we have the exact row-replacement identity
\[
 \mathcal C_j=\mathcal C_0+e_jc_j^{\top}.                         \tag{3}
\]
Relabel only the latent nodes and put
\[
 d=p,
 \qquad
 G_{\mathrm{Sq}}=Q^{\top},
 \qquad
 H_{\mathrm{Sq}}=G_{\mathrm{Sq}}^{\dagger}=Q,
 \qquad
 \mathcal B_c^{\mathrm{Sq}}=Q\mathcal C_cQ^{\top}
 \quad(c\in C).                                                  \tag{4}
\]
Equation (1) implies that \(\mathcal B_0^{\mathrm{Sq}}\) is upper
triangular in the convention of Squires et al.: an original edge
\(u\to v\) contributes an entry in position
\((\rho(v),\rho(u))\), and \(\rho(v)<\rho(u)\).  Its diagonal entries
are the positive numbers \(\omega_j^{-1/2}\), merely permuted.  Thus it
is the normalized mechanism matrix of the reverse-topologically relabeled
baseline DAG.  The same argument, with the target row replaced, shows
that every \(\mathcal B_j^{\mathrm{Sq}}\) is nonsingular and represents
the corresponding perfect intervention.

Multiplying (3) on both sides by the permutation matrices gives
\[
 \begin{aligned}
 \mathcal B_j^{\mathrm{Sq}}
 &=\mathcal B_0^{\mathrm{Sq}}+(Qe_j)(Qc_j)^{\top}\\
 &=\mathcal B_0^{\mathrm{Sq}}+e_{\rho(j)}
 \left(\tau_j^{-1/2}e_{\rho(j)}
       -(\mathcal B_0^{\mathrm{Sq}})^{\top}e_{\rho(j)}\right)^{\top}.
                                                                    \tag{5}
 \end{aligned}
\]
Thus context \(j\) targets latent node \(\rho(j)\), and (5) is exactly
the paper's single-row perfect-intervention form with positive intervention
scale \(\tau_j^{-1/2}\).

We next verify the two conditions in Assumption 1 of Squires
et al.~\cite[Assumption 1]{pmlr-v202-squires23a} that
are not already supplied by the normalized acyclic representation.  First,
if \(j\) is a nonroot, choose \(i\in\operatorname{Pa}(j)\).  By the
definition of \(B\), the baseline row \(j\) has the nonzero parent entry
\(-\omega_j^{-1/2}B_{ji}\), whereas the intervention row is supported
only at coordinate \(j\).  After applying \(Q\),
\((\mathcal B_0^{\mathrm{Sq}})^{\top}e_{\rho(j)}\) therefore has a
nonzero entry at \(\rho(i)\), while
\[
 (\mathcal B_j^{\mathrm{Sq}})^{\top}e_{\rho(j)}
 =\tau_j^{-1/2}e_{\rho(j)}.
\]
The two vectors cannot be proportional.  If \(j\) is a root, the
baseline vector is \(\omega_j^{-1/2}e_{\rho(j)}\), so proportionality is
possible and is expressly allowed by Assumption 1(b).  Second,
\(G_{\mathrm{Sq}}\) is full rank and its Moore--Penrose inverse is
\(H_{\mathrm{Sq}}=Q\).  Each row of \(Q\) has exactly one entry equal to
one and all other entries zero.  Its largest absolute entry is therefore
one, with no tie, which verifies the row normalization in Assumption
1(c).  Consequently (4)--(5), together with the centered unit-variance
noise representation
\(Z_c=(\mathcal B_c^{\mathrm{Sq}})^{-1}\varepsilon\),
\(\varepsilon\sim\mathcal N(0,I_p)\), satisfies the published model
assumptions.

It remains to check that this latent relabeling changes none of the
observable precision matrices.  By \(\texttt{def:model-class}\), direct
inversion of the baseline and surgery covariance factorizations gives
\[
 K_0=A^{\top}\Omega^{-1}A=\mathcal C_0^{\top}\mathcal C_0,
 \qquad
 K_j=(A^{(j)})^{\top}(\Omega^{(j)})^{-1}A^{(j)}
     =\mathcal C_j^{\top}\mathcal C_j.                            \tag{6}
\]
Using (4), orthogonality of \(Q\), and the source's observed-precision
formula gives, for every \(c\in C\),
\[
 \begin{aligned}
 \Theta_c
 &=H_{\mathrm{Sq}}^{\top}
   (\mathcal B_c^{\mathrm{Sq}})^{\top}
   \mathcal B_c^{\mathrm{Sq}}H_{\mathrm{Sq}}\\
 &=Q^{\top}(Q\mathcal C_c^{\top}Q^{\top})
            (Q\mathcal C_cQ^{\top})Q\\
 &=\mathcal C_c^{\top}\mathcal C_c
 =K_c.                                                            \tag{7}
 \end{aligned}
\]
In particular, the factors \(H_{\mathrm{Sq}}\) have not been dropped.
No context label has been permuted: label \(j\) still names context
\(j\), although its latent target is \(\rho(j)\).  Accordingly, if
\(\sigma=(0\;1)\) is used in the rank-only argument, it denotes only the
transposition of two component labels in \(C\), and it is distinct from
the latent-node bijection \(\rho\).  Finally, (7) yields
\[
 \operatorname{rank}(\Theta_c-\Theta_d)
 =\operatorname{rank}(K_c-K_d)
 \quad(c,d\in C).
\]
Because the context labels are unchanged, every entry of the rank matrix,
the literal Proposition-10 score, and the symmetric score is preserved.
\end{proof}
\par\smallskip\noindent \textit{Justification.} The reverse-topological latent relabeling verifies the published order, generic-row, and row-normalization conditions and proves \(\Theta_c=K_c\) without dropping the observation-map factors. \textit{Gap.} The former identity-mixing witness did not satisfy the published latent ordering for arbitrary coordinate labels. \textit{Consumer.} The full-domain negative transfer uses this lemma once; downstream proofs need not recreate or alter the embedding.

\begin{lemma}[lem:rank-witness-calculations]\label{lem:rank-witness-calculations}
The two chain specializations in \(\texttt{def:two-node-counterfamilies}\), the fork in \(\texttt{def:three-node-counterfamily}\), and the strict one-root-star region in \(\texttt{def:star-counterfamily}\) are legal members of their respective classes \(\mathcal M_p\).  The first chain has literal scores \((3,2,2)\) and symmetric scores \((3,3,4)\); the second has symmetric scores \((3,2,3)\).  The fork has literal scores \((5,4,5,5)\).  More generally, throughout the star region, the published literal score of the baseline is \(2p-1\), whereas that of the root intervention is \(2p-2\).  The coordinate certificate accepts every one of these tuples and returns its true baseline and semantic explanation.
\end{lemma}
\begin{proof}
For the chain \(1\to2\), write \(B_{21}=b\neq0\), \(\omega=(u,v)\), and \(\tau=(t,s)\).  Baseline factorization and the surgery identity give
\[
 K_0=
 \begin{pmatrix}
 u^{-1}+b^2v^{-1}&-bv^{-1}\\
 -bv^{-1}&v^{-1}
 \end{pmatrix},\qquad
 K_1=
 \begin{pmatrix}
 t^{-1}+b^2v^{-1}&-bv^{-1}\\
 -bv^{-1}&v^{-1}
 \end{pmatrix},\qquad
 K_2=
 \begin{pmatrix}
 u^{-1}&0\\0&s^{-1}
 \end{pmatrix}.                                             \tag{27}
\]
All five scalar parameters in both displayed specializations are positive, \(b\neq0\), and the root intervention is effective because \(t\neq u\).  Target \(2\) is nonroot because \(b\neq0\), so its intervention is effective whether or not \(s=v\).  Hence both are in \(\mathcal M_2\) by \(\texttt{def:model-class}\).

At \((b,u,v,t,s)=(2,1,1,3,2)\), (27) becomes
\[
 K_0=\begin{pmatrix}5&-2\\-2&1\end{pmatrix},\qquad
 K_1=\begin{pmatrix}13/3&-2\\-2&1\end{pmatrix},\qquad
 K_2=\begin{pmatrix}1&0\\0&1/2\end{pmatrix}.
\]
Here \(K_0-K_1\) is nonzero diagonal of rank one.  Moreover,
\[
 \det(K_0-K_2)=-2\neq0,
 \qquad
 \det(K_1-K_2)=-7/3\neq0.
\]
Thus
\[
 R=\begin{pmatrix}0&1&2\\1&0&2\\2&2&0\end{pmatrix}.       \tag{28}
\]
The literal convention in Proposition 10 sums columns \(1,2\), giving \((3,2,2)\); summing instead over every component other than the scored component gives \((3,3,4)\).

At \((b,u,v,t,s)=(1,1,1,1/2,2)\),
\[
 K_0=\begin{pmatrix}2&-1\\-1&1\end{pmatrix},\qquad
 K_1=\begin{pmatrix}3&-1\\-1&1\end{pmatrix},\qquad
 K_2=\begin{pmatrix}1&0\\0&1/2\end{pmatrix}.
\]
The first two pairwise differences have ranks one and two respectively, while \(K_1-K_2\neq0\) and
\[
 \det(K_1-K_2)=2(1/2)-1=0.
\]
It therefore has rank one, so
\[
 R=\begin{pmatrix}0&1&2\\1&0&1\\2&1&0\end{pmatrix},
\]
whose symmetric scores are \((3,2,3)\).

For the fork \(1\to2,1\to3\), all baseline and intervention variances are positive, the root variance changes from \(1\) to \(3\), and each other target is nonroot.  It is therefore legal.  The same factorization gives
\[
 \begin{aligned}
 K_0&=\begin{pmatrix}14&-2&-3\\-2&1&0\\-3&0&1\end{pmatrix},&
 K_1&=\begin{pmatrix}40/3&-2&-3\\-2&1&0\\-3&0&1\end{pmatrix},\\
 K_2&=\begin{pmatrix}10&0&-3\\0&1/4&0\\-3&0&1\end{pmatrix},&
 K_3&=\begin{pmatrix}5&-2&0\\-2&1&0\\0&0&1/5\end{pmatrix}.
 \end{aligned}                                               \tag{29}
\]
The baseline--root difference has rank one.  The principal \(2\)-by-\(2\) determinants on coordinates \(\{1,2\}\) for \(K_0-K_2\) and \(K_1-K_2\) are \(-1\) and \(-3/2\); those on \(\{1,3\}\) for \(K_0-K_3\) and \(K_1-K_3\) are \(-9/5\) and \(-7/3\).  Each corresponding difference is supported on those two coordinates and hence has rank two.  Finally,
\[
 \det(K_2-K_3)=11/20\neq0,
\]
so that difference has rank three.  Therefore
\[
 R=\begin{pmatrix}
 0&1&2&2\\1&0&2&2\\2&2&0&3\\2&2&3&0
 \end{pmatrix},
\]
and summation over the published columns \(1,2,3\) gives \((5,4,5,5)\).

It remains to verify the claimed open family.  Fix the star of \(\texttt{def:star-counterfamily}\), let
\[
 d=\tau_1^{-1}-\omega_1^{-1}\neq0,
\]
and take a leaf \(j\ge2\).  By \(\texttt{lem:surgery-identity}\),
\[
 K_1-K_0=d e_1e_1^{\top},
 \qquad
 (K_j-K_0)_{\{1,j\},\{1,j\}}
 =\begin{pmatrix}
 -b_j^2/\omega_j&b_j/\omega_j\\
 b_j/\omega_j&\tau_j^{-1}-\omega_j^{-1}
 \end{pmatrix}.                                               \tag{30}
\]
The determinant of the second matrix in (30) is
\[
 -\frac{b_j^2}{\omega_j\tau_j}<0,
\]
so \(\operatorname{rank}(K_j-K_0)=2\).  Subtracting the root difference changes only entry \((1,1)\), and direct expansion yields
\[
 \det\bigl((K_j-K_1)_{\{1,j\},\{1,j\}}\bigr)
 =\frac{d(\tau_j-\omega_j)-b_j^2}{\omega_j\tau_j}\neq0,       \tag{31}
\]
where the last inequality is the strict star condition.  Hence \(\operatorname{rank}(K_j-K_1)=2\).  The region is nonempty: taking every \(\omega_\ell=1\), \(\tau_1=2\), and \(\tau_j=b_j=1\) for all leaves gives \(d=-1/2\) and makes every numerator in (31) equal to \(-1\).  Positivity, \(d\neq0\), \(b_j\neq0\), and the finitely many strict inequalities persist on an open neighborhood.  The changed root variance and the nonzero parent coefficient of each leaf verify effectiveness, so the neighborhood lies in \(\mathcal M_p\).

For every point of this region,
\[
 R_{01}=1,
 \qquad R_{0j}=R_{1j}=2\quad(2\le j\le p).
\]
Thus the published score \(q_c=\sum_{\ell=1}^pR_{c\ell}\) satisfies
\[
 q_0=1+2(p-1)=2p-1,
 \qquad
 q_1=2(p-1)=2p-2.                                           \tag{32}
\]
Finally, every tuple just considered is a member of \(\mathcal M_p\).  Completeness and uniqueness in \(\texttt{thm:complete-certificate}\) therefore imply that \(\mathsf{Cert}\) accepts it and returns its true baseline and semantic explanation.
\end{proof}

\begin{proposition}[prop:squires-rank-failure-full-domain]\label{prop:squires-rank-failure-full-domain}
The chain, fork, and open star witnesses falsify the observational-context rank-selector conclusion of Squires et al., Appendix F.2, Proposition 10 on its full published model class, not merely on \(\mathcal M_p\): under the normalized-mechanism embedding their context precisions, rank matrices, and scores are unchanged, while the true observational context is not a literal-score minimizer in the first chain, the fork, and the entire star region.  The natural symmetric score also fails, by a tie in the first chain and a strict error in the second.  This is a negative scope transfer only; the positive coordinate certificate is proved on the observed-coordinate fully covered class \(\mathcal M_p\) and does not assert recovery under latent mixing.
\end{proposition}
\begin{proof}
By \(\texttt{lem:squires-class-embedding}\), with its reverse-topological
latent relabeling, every legal witness in \(\mathcal M_p\) satisfies the
full published assumptions of Squires
et al.~\cite[Assumption 1 and Appendix F.2, Proposition 10]{pmlr-v202-squires23a}
and has
\(\Theta_c=K_c\) at the same context label \(c\).  Therefore all
pairwise precision-difference ranks and both score conventions are
unchanged.  The embedding is invoked here as a proved lemma; no second
mixing construction is being assumed.

The exact calculations in \(\texttt{lem:rank-witness-calculations}\)
give the following failures.  In the first chain, the observational
context has literal score \(3\), while both intervention contexts have
score \(2\).  In the fork, the observational context has literal score
\(5\), while the root-intervention context has score \(4\).  Throughout
the nonempty open star region, the corresponding scores are
\[
 2p-1\quad\text{and}\quad 2p-2,
\]
respectively.  In every case the true observational context is not a
literal-score minimizer.  These are points in the full published class by
the embedding, so they falsify the conclusion of Appendix F.2,
Proposition 10 on that class.  The same calculation gives symmetric scores
\((3,3,4)\) in the first chain, which destroys uniqueness, and
\((3,2,3)\) in the second chain, which strictly selects the wrong
context.

This transfer is one-way.  The counterexamples belong to the full class
because the embedded subclass is contained in it.  By contrast,
\(\texttt{thm:complete-certificate}\) proves acceptance, unique recovery,
and rejection traces only for tuples having an explanation in
\(\mathcal M_p\).  Thus the negative rank conclusion holds on the full
published domain by restriction, while no positive certificate or
recovery claim is made for arbitrary latent mixing.
\end{proof}
\par\smallskip\noindent \textit{Justification.} The repaired embedding places each legal chain, fork, and open star witness in the full published class with its context labels, precision matrix, and rank scores unchanged. \textit{Gap.} A counterexample on \(\mathcal M_p\) refutes the larger-class claim only after the ordering, intervention-genericity, observation normalization, and precision-preservation conditions are all verified. \textit{Consumer.} Readers obtain a valid negative scope transfer while retaining the positive coordinate certificate only on \(\mathcal M_p\).

\section*{Technical internal limitation (diagnostic only)}

The exact certificate is a real-arithmetic population procedure. It does not test approximate equalities, certify rational bit complexity, recover individual latent regime labels for pooled observations, or provide an end-to-end component estimator. The causaldag documentation for unknown-target IGSP expects separate observational and interventional sufficient statistics and regime metadata. On an accepted exact tuple, the certificate can supply the semantic baseline and target labels that such a workflow needs, and it recovers the DAG directly within its narrower model. It cannot by itself manufacture segregated samples or the empirical conditional-independence and invariance testers required by that software interface; it is therefore a semantic validation and initialization layer, not a drop-in statistical replacement.

\section{Honest scope}

The population identification, exact membership, and deterministic stability results are confined to the observed-coordinate, fully covered perfect-surgery class \(\mathcal M_p\). The negative comparison is broader: a reverse-topological latent-node bijection \(\rho\), its permutation matrix \(Q\), \(G_{\mathrm{Sq}}=Q^{\top}\), \(H_{\mathrm{Sq}}=Q\), and conjugated normalized mechanism matrices embed \(\mathcal M_p\) into the full class of Squires et al. without changing context labels or context precisions. Hence the rank-selector counterexamples and the equivariant rank-only impossibility hold on that full class by restriction. No positive recovery or certificate under arbitrary latent mixing is claimed. The exact certificate can semantically validate and initialize the documented causaldag unknown-target IGSP workflow by supplying baseline and target labels for an accepted population tuple, but it does not create segregated samples, regime metadata, or empirical testers. The growing-dimensional component operator-error tail remains open. The restriction \(p\ge2\) is necessary because at \(p=1\) the two distinct variances admit both baseline orientations. Missing interventions, soft or multitarget interventions, cycles, singular components, and latent coordinate mixing remain outside the positive theorem.

\begin{thebibliography}{99}

  \bibitem{pmlr-v202-squires23a} Squires, C., Seigal, A., Bhate, S. S., and Uhler, C. (2023). Linear Causal Disentanglement via Interventions. Proceedings of the 40th International Conference on Machine Learning, PMLR 202, 32540--32560. Theorem 1, Proposition 1, Appendix F.2 Proposition 10, and Appendix G.

  \bibitem{kumarShiragurUhler2024mixtures} Kumar, A., Shiragur, K., and Uhler, C. (2024). Learning Mixtures of Unknown Causal Interventions. Advances in Neural Information Processing Systems 37. Theorem 4.1, Corollary 4.1.1, and Section 8.

  \bibitem{kumarSinha2021disentangling} Kumar, A., and Sinha, G. (2021). Disentangling Mixtures of Unknown Causal Interventions. Proceedings of the Thirty-Seventh Conference on Uncertainty in Artificial Intelligence, PMLR 161, 2093--2102. Definition 3.1, Assumptions 3.1 and 3.2, and Theorem 3.1.

  \bibitem{variciEtAl2021cite} Varici, B., Shanmugam, K., Sattigeri, P., and Tajer, A. (2021). Scalable Intervention Target Estimation in Linear Models. Advances in Neural Information Processing Systems 34. Equation (6) and Theorems 1 and 3.

  \bibitem{squiresWangUhler2020utigsp} Squires, C., Wang, Y., and Uhler, C. (2020). Permutation-Based Causal Structure Learning with Unknown Intervention Targets. Proceedings of the 36th Conference on Uncertainty in Artificial Intelligence, PMLR 124. Algorithm 1 and Theorem 1.

  \bibitem{gamellaEtAl2022gnies} Gamella, J. L., Taeb, A., Heinze-Deml, C., and Buehlmann, P. (2022). Characterization and Greedy Learning of Gaussian Structural Causal Models under Unknown Interventions. arXiv:2211.14897. Theorem 6 in the 2022 manuscript, renumbered as Theorem 1 in arXiv version 4; Definition 1 and Section 6.

  \bibitem{andersonEtAl2026dimension} Anderson, P., Bafna, M., Buhai, R.-D., Kothari, P. K., and Steurer, D. (2026). Dimension Reduction via Sum-of-Squares and Improved Clustering Algorithms for Non-Spherical Mixtures. arXiv:2411.12438, version 2. Theorem 1.1 and Section 2.1.

  \bibitem{koltchinskiiLounici2017covariance} Koltchinskii, V., and Lounici, K. (2017). Concentration Inequalities and Moment Bounds for Sample Covariance Operators. Bernoulli 23, 110--133. Theorems 4 and 6.

  \bibitem{leyesCarrenoMeroniSeigal2025} Leyes Carreno, P., Meroni, C., and Seigal, A. (2025). Linear Causal Disentanglement via Higher-Order Cumulants. La Matematica. Definition 1.5 for identifiability terminology, Theorem 1.6 for perfect-intervention recovery under a known observational context and non-Gaussian errors, and Section 4.

  \bibitem{causaldag2020utigsp} Uhler Lab (2020). UT-IGSP documentation in the causaldag package. The documented example constructs separate observational and interventional sufficient statistics and calls the unknown-target routine with regime metadata.

\end{thebibliography}

\end{document}